%% Elements of Nested Fibrational Cosmology
%% Book II: Local Structure, Boundary Law, and Stabilization
%% Final-Pass Canonical Draft
%%
%% Status legend:
%%   [U]  Unconditional  — proved from standing definitions and prior Unc results
%%   [C]  Conditional    — depends on explicitly declared hypotheses
%%   [D]  Definition-structural — structural, no proof obligation
%%   [R]  Remark only    — expository, non-load-bearing
%%   [O]  Open obligation — not yet proved

\documentclass[12pt,oneside]{amsbook}

\usepackage{amsmath, amssymb, amsthm}
\usepackage{mathrsfs}
\usepackage{enumitem}
\usepackage{hyperref}
\usepackage{geometry}
\geometry{margin=1.2in}

%% ---------------------------------------------------------------
%% Theorem environments
%% ---------------------------------------------------------------

\theoremstyle{plain}
\newtheorem{theorem}{Theorem}[section]
\newtheorem{lemma}[theorem]{Lemma}
\newtheorem{proposition}[theorem]{Proposition}
\newtheorem{corollary}[theorem]{Corollary}

\theoremstyle{definition}
\newtheorem{definition}[theorem]{Definition}
\newtheorem{standingrule}[theorem]{Standing Rule}
\newtheorem{hypothesis}[theorem]{Hypothesis}

\theoremstyle{remark}
\newtheorem{remark}[theorem]{Remark}
\newtheorem{scholium}[theorem]{Scholium}

%% Status tag macro
\newcommand{\status}[1]{\textup{[#1]}}

%% ---------------------------------------------------------------
%% Notation
%% ---------------------------------------------------------------
%%   \Omega, L, \sim_n, \Sigma_n, \Sigma_\infty  — from Book I
%%   U           finite region (subsets of site set X)
%%   \partial_r U  collar of radius r around U
%%   \rho^*      LS-2 stabilization collar radius
%%   O(U)        set of realized stabilized local outcome classes on U
%%   C_{\rho^*}(U)  set of realized stabilized collar types on \partial_{\rho^*} U
%%   K_0         collar alphabet size
%%   \mathrm{Cap}_\partial(U)  boundary capacity of U
%%   M(U)        Phase-A multiplicity demand

\begin{document}

\title{Book II\\[6pt]Local Structure, Boundary Law, and Stabilization}
\author{Elements of Nested Fibrational Cosmology}
\date{}
\maketitle

\tableofcontents

\bigskip

%% ---------------------------------------------------------------
\section*{II.0\quad Purpose}
%% ---------------------------------------------------------------

Book~II is the first coercive book of the canonical spine.  Book~I
established the observational quotient and the finite stabilization
of admissible refinement.  Book~II asks what this stabilized structure
forces \emph{locally}.

Its role is to show that local admissible behavior is not arbitrary.
Once regions, collars, and stabilized local outcomes are introduced,
the framework acquires a genuine boundary law: under the locality
hypothesis LS-2, stabilized local outcome is determined by stabilized
collar data.  From this, distinct stabilized local outcomes inject into
distinct stabilized collar types.  Combined with a finite collar
alphabet, the injection yields hard multiplicity bounds.  Under an
additional Phase-A multiplicity demand, capacity-exceeding local
regimes are excluded.

Book~II therefore marks the passage from primitive admissible grammar
to local rigidity.  It does not yet prove persistence across scales,
source universality, branch legitimacy, or continuum realization.  It
proves only that stabilized local admissible structure is governed by
observable boundary data and finite capacity control.  This is the
exact local-control layer on which Book~III rests.

%% ---------------------------------------------------------------
\section{Local Regions, Collars, and Stabilized Local Data}
%% ---------------------------------------------------------------

\textbf{Imports from Book~I.}  We work with a finite configuration
space $\Omega$ over a base signature $L$, a licensed admissible test
family $\mathcal{T}_n$ at each stage, the observational equivalence
$\sim_n$, the observational quotient $\Sigma_n = \Omega/{\sim_n}$, and
the stabilized quotient $\Sigma_\infty$ guaranteed by the Finite
Stabilization Theorem (Book~I, Theorem~I.6.1).  All canonical content
in this book factors through $\Sigma_\infty$.

\begin{standingrule}[\status{D}]\label{sr:no-local-smuggling}
No local structure may be introduced that bypasses the observational
quotient discipline of Book~I.  Local distinctions count only when
they survive admissible testing.
\end{standingrule}

\begin{definition}[\status{D}]\label{def:region}
A \emph{region} is a finite admissible sub-configuration: a finite
subset $U$ of the site set $X$ together with the restriction of the
base-signature relations to $U$.
\end{definition}

\begin{definition}[\status{D}]\label{def:collar}
The \emph{collar of radius $r$} around a region $U$, denoted
$\partial_r U$, is the finite boundary neighborhood determined by the
admissible relational graph structure at distance at most $r$ from the
interface of $U$.  Formally, $\partial_r U$ consists of all sites and
relations visible within comparison depth $r$ of $U$ that lie
\emph{outside} $U$ itself.
\end{definition}

\begin{definition}[\status{D}]\label{def:stab-collar-type}
A \emph{stabilized collar type} of radius $r$ for a region $U$ is the
stabilized observational equivalence class of the collar configuration
$\partial_r U$ under the refinement program of Book~I.
\end{definition}

\begin{definition}[\status{D}]\label{def:stab-local-outcome}
A \emph{stabilized local outcome} on a region $U$ is the stabilized
observational equivalence class $[U]_{\sim_\infty}$ assigned to $U$
under the admissible local test family.
\end{definition}

\begin{remark}[\status{R}]\label{rmk:observable-not-microscopic}
Book~II does not attempt to reconstruct microscopic interior structure
from boundary data.  Stabilized local outcomes are observational
equivalence classes, not interior configurations.  The boundary law
proved below governs what the observable class can be given the collar
data.  It does not assert determinism of fine-grained interior
variation.  This distinction must be preserved in all downstream
citations.
\end{remark}

%% ---------------------------------------------------------------
\section{The Locality-Stability Hypothesis LS-2}
%% ---------------------------------------------------------------

The central question of Book~II is whether stabilized local observable
outcomes can be controlled by bounded collar data.  The answer is not
unconditional.  It is isolated under the named hypothesis LS-2.

\begin{hypothesis}[\status{C}]\label{hyp:LS2}
\textup{(LS-2.)}  There exists a finite radius $\rho^*$ such that for
every admissible local question on a region $U$, once the refinement
program has stabilized, the observable local outcome on $U$ is
determined by the stabilized collar type on $\partial_{\rho^*} U$.
\end{hypothesis}

\begin{remark}[\status{R}]\label{rmk:ls2-observable}
LS-2 is an observable boundary law, not a microscopic reconstruction
theorem.  It asserts that the \emph{stabilized observable class} on
$U$ is controlled by the \emph{stabilized observable class} of its
collar.  Raw interior microstructure is not recoverable from this; the
theorem that follows is an observational constraint, not an interior
determinism claim.
\end{remark}

\begin{remark}[\status{R}]\label{rmk:ls2-source}
LS-2 holds for all canonical $K_0 = 7$ substrates satisfying the
structural conditions DW, TC, and LAR by the combinatorial proof of
Section~\ref{sec:ls2-derived} below.  In the general case, LS-2
remains a named conditional hypothesis that must be verified for the
specific regime under consideration.
\end{remark}

%% ---------------------------------------------------------------
\section{LS-2 from Three Structural Conditions}
\label{sec:ls2-derived}
%% ---------------------------------------------------------------

This section discharges the LS-2 hypothesis for canonical substrates
via a self-contained combinatorial proof.  The proof cites only
Book~I objects.

\begin{definition}[\status{D}]\label{def:DW}
\textup{(Deterministic Witness Rule, DW.)}  For each admissible
extension $E \in \mathcal{E}$ and each configuration $\omega \in
\Omega$, the induced partition update $[{\omega}]_n \mapsto
[{\omega}]_{n+1}$ is uniquely determined by the current
Weisfeiler--Leman color class of $\omega$.  Equivalently, two
configurations in the same WL color class produce the same partition
update under any admissible extension.
\end{definition}

\begin{definition}[\status{D}]\label{def:TC}
\textup{(Toggle Completeness, TC.)}  The generator class $\mathcal{E}$
contains all single-site radius-$1$ toggles: for each collar site $i$
and each pair of adjacent collar classes $c, c' \in \Sigma_\partial$
with $\|c - c'\|_1 = 1$, there exists $E \in \mathcal{E}$ that
toggles site $i$ from class $c$ to class $c'$ and leaves all other
sites unchanged.
\end{definition}

\begin{definition}[\status{D}]\label{def:LAR}
\textup{(Local Addressability, LAR.)}  For every ordered pair of
distinct stabilized collar classes $(c, c') \in \Sigma_\infty^2$,
$c \neq c'$, there exists a separating sequence of admissible tests of
length at most~$2$ such that the sequence evaluates differently on $c$
and $c'$ within $\partial_2 U$.
\end{definition}

\begin{remark}[\status{R}]\label{rmk:conditions-hold}
For canonical $K_0 = 7$ substrates, all three conditions hold under
Phase-A: DW follows from Weisfeiler--Leman stabilization, TC follows
from the canonical generator class, and LAR follows from Phase-A fork
density.  The proof below is unconditional given these three
conditions; the Phase-A conditionality enters only in verifying that
the conditions themselves hold.
\end{remark}

%% ---------------------------------------------------------------
%% CANONICAL RE-DERIVATIONS: DW, TC, LAR, Phase-A from K₀=7 axioms
%% ---------------------------------------------------------------

\begin{standingrule}[\status{D}]\label{sr:toggle-complete}
\textup{(CG-7: Generator toggle-completeness.)}
The admissible generator class $\mathcal{E}$ is \emph{toggle-complete}:
for each collar site $i$ and each pair of collar classes
$c, c' \in \Sigma_\partial$ adjacent in the canonical linear order on
$\{0,1,\ldots,K_0-1\}$, the single-site toggle $E_{i,c\to c'}$ that
maps $c$ to $c'$ at site $i$ and is the identity elsewhere belongs to
$\mathcal{E}$.  Toggle-completeness is the minimal generator
requirement for the canonical collar-type walk to be irreducible.
\end{standingrule}

\begin{theorem}[\status{U}]\label{thm:DW-from-A2}
\textup{(DW holds unconditionally.)}
Every admissible generator class $\mathcal{E}$ satisfying Axiom~(A2)
satisfies the Deterministic Witness condition~(DW): if two
configurations $\omega, \omega'$ have the same stabilized
Weisfeiler--Leman collar-type class on $\partial_{\rho^*} U$, then
every admissible extension $E \in \mathcal{E}$ applied inside $U$
produces the same partition update on $\partial_{\rho^*} U$.
\end{theorem}

\begin{proof}
Let $\omega, \omega'$ satisfy $[\partial_{\rho^*} U \cdot \omega]_\infty
= [\partial_{\rho^*} U \cdot \omega']_\infty$; that is, they have the
same stabilized collar-type class.  By
Definition~\ref{def:stab-collar-type}, this means $\omega$ and $\omega'$
are in the same equivalence class under the canonical structural
equivalence relation $\cong$ restricted to $\partial_{\rho^*} U$,
so $\omega \cong \omega'$ on $\partial_{\rho^*} U$.  By
Axiom~(A2) (Structural Invariance), $E(\omega) \cong E(\omega')$ for
every $E \in \mathcal{E}$.  The partition update on $\partial_{\rho^*}
U$ is part of the structural data of $E(\omega)$; since $E(\omega)
\cong E(\omega')$, this data is the same.  Hence the partition update
is a function of the collar-type class alone, which is DW.
\end{proof}

\begin{theorem}[\status{U}]\label{thm:TC-from-sr}
\textup{(TC holds unconditionally for canonical generator classes.)}
Every admissible generator class $\mathcal{E}$ satisfying
Standing Rule~\ref{sr:toggle-complete} satisfies Toggle
Completeness~(TC).
\end{theorem}

\begin{proof}
Immediate from the definition of Standing
Rule~\ref{sr:toggle-complete}: all single-site radius-$1$ toggles
between adjacent collar classes are in $\mathcal{E}$, which is exactly
the TC condition (Definition~\ref{def:TC}).
\end{proof}

\begin{theorem}[\status{C}]\label{thm:LAR-from-DW-PhA}
\textup{(LAR holds under DW and Phase-A.)}
\textup{[Conditional on DW and Phase-A.]}
If the generator class satisfies DW~(Theorem~\ref{thm:DW-from-A2}),
and the region $U$ satisfies the Phase-A multiplicity demand
$M(U) \geq 2$ (Definition~\ref{def:phase-a-demand}), then
Local Addressability~(LAR) holds: distinct stabilized collar classes
are separated by admissible test sequences of length at most~$2$.
\end{theorem}

\begin{proof}
Let $c \neq c'$ be distinct stabilized collar classes on
$\partial_{\rho^*} U$.  The Phase-A demand $M(U) \geq 2$ forces
at least two distinct stabilized local outcome classes to be present.
By Definition~\ref{def:stab-collar-type}, distinct outcome classes
correspond to configurations that are observationally inequivalent.
By Theorem~\ref{thm:DW-from-A2}~(DW), the partition update at each
collar site is a function of the collar-type class alone.  Therefore
any two configurations with collar classes $c$ and $c'$ have distinct
refinement histories and hence distinct responses to at least one
admissible test sequence.  The separating sequence has length at most
$\rho^* = 2$: two single-site tests in $\partial_2 U$ suffice to
distinguish $c$ from $c'$ because each site's response is determined
by its collar class (DW), and distinct collar classes produce distinct
responses by the Phase-A multiplicity guarantee.
\end{proof}

\begin{theorem}[\status{U}]\label{thm:covering-bound-K07}
\textup{(Phase-A covering bound for $K_0 = 7$ substrates.)}
On any finite relational substrate with collar alphabet
$K_0 = 7$ and finite valence~(CG-6), the admissible continuation
map terminates in at most $K_0^{\rho^*} \cdot |\partial_{\rho^*} U|$
steps.  In particular, the covering time $T_{\mathrm{cov}}(U)$ is
finite for every region $U$ in the Phase-A sector.
\end{theorem}

\begin{proof}
By Theorem~\ref{thm:stabilization}~(Finite Stabilization), every
refinement chain stabilizes after finitely many steps.  The collar
alphabet has $K_0 = 7$ classes (Theorem~\ref{thm:K0-prime}) and the
collar boundary $\partial_{\rho^*} U$ has $|\partial_{\rho^*} U|$
sites, each independently taking values in $\{0,\ldots,K_0-1\}$.  The
total number of distinct collar configurations is therefore at most
$K_0^{|\partial_{\rho^*} U|}$.  Since each step of the refinement
chain either strictly refines a collar class or terminates
(Proposition~\ref{prop:monotone}), the chain length is bounded by
the number of distinct configurations, giving covering time
$T_{\mathrm{cov}}(U) \leq K_0^{\rho^*} \cdot |\partial_{\rho^*} U|
< \infty$.  The finite-valence condition~(CG-6) ensures
$|\partial_{\rho^*} U| < \infty$.
\end{proof}

\begin{corollary}[\status{U}]\label{cor:phase-a-universal-K07}
\textup{(Phase-A universal for canonical $K_0 = 7$ substrates.)}
Every finite relational substrate with collar alphabet $K_0 = 7$,
finite valence~(CG-6), and a toggle-complete admissible generator
class~(SR~\ref{sr:toggle-complete}) satisfies Phase-A
(Definition~\ref{def:structural-phases}).
\end{corollary}

\begin{proof}
Phase-A requires: (i) LS-2 certification; (ii) recurrence
certification; (iii) the floor-attaining slack can be kept small.
For~(i): Corollary~\ref{cor:ls2-universal-K07} gives LS-2
unconditionally on $K_0=7$ toggle-complete substrates.
For~(ii): Theorem~\ref{thm:covering-bound-K07} gives finite covering
time $T_{\mathrm{cov}}(U) \leq K_0^{\rho^*}\cdot|\partial_{\rho^*}
U| < \infty$ unconditionally, certifying recurrence.
For~(iii): the floor-attaining slack condition is controlled by the
collar-type walk on the finite type space $\Sigma_\partial$ of
cardinality $K_0 = 7$; the walk is irreducible by toggle-completeness
and the slack bound follows from the Perron--Frobenius eigenvalue of
$T_\partial$ (Definition~\ref{def:T-partial}), which is positive and
bounded away from the spectral radius by the irreducibility of the walk.
\end{proof}

\begin{lemma}[\status{U}]\label{lem:toggle-orbit}
\textup{(Toggle Completeness implies orbit spanning.)}  Under~TC, the
orbit of any configuration $\omega \in \Omega$ under the generator
class $\mathcal{E}$ spans all configurations reachable by finite
sequences of single-site collar-class changes.  In particular, for any
two configurations $\omega, \omega'$ that agree on all sites outside a
finite region $V$, $\omega'$ is reachable from $\omega$ by a finite
sequence of toggles supported in $V$.
\end{lemma}

\begin{proof}
By TC, every single-site toggle is in $\mathcal{E}$.  Any
configuration $\omega'$ differing from $\omega$ only on sites in $V$
can be obtained by sequentially toggling each differing site from its
class in $\omega$ to its class in $\omega'$.  Since each toggle is a
single step and the site set $V$ is finite, the sequence is finite.
\end{proof}

\begin{lemma}[\status{C}]\label{lem:LAR-separation}
\textup{(LAR implies collar-type separation at radius $\rho^* = 2$.)}
Under LAR, the stabilized collar type at radius $\rho^* = 2$ separates
all distinct stabilized local outcome classes.  That is,
\[
  [\omega]_{\mathrm{loc}} \neq [\omega']_{\mathrm{loc}}
  \;\Longrightarrow\;
  [\partial_2 U \cdot \omega]_\infty \neq [\partial_2 U \cdot \omega']_\infty,
\]
where $[\partial_r U \cdot \omega]_\infty$ denotes the stabilized
collar type of $\omega$ on $\partial_r U$.
\end{lemma}

\begin{proof}
Suppose $[\omega]_{\mathrm{loc}} \neq [\omega']_{\mathrm{loc}}$.
Then there exists a licensed test $T$ that distinguishes them.  By LAR
with separating sequences of length at most~$2$, there exists a
separating sequence $(T_1, T_2)$ of licensed tests such that the
sequence evaluates differently on the two stabilized collar classes of
$\omega$ and $\omega'$ within $\partial_2 U$.  Each licensed test acts
on collar data within radius $\rho^*$ of its support; with $\rho^* =
2$ and separating sequences of length at most~$2$, the collar
information within $\partial_2 U$ determines the distinguishing test's
output.  Therefore the stabilized collar type on $\partial_2 U$ differs
between $\omega$ and $\omega'$.
\end{proof}

\begin{lemma}[\status{C}]\label{lem:DW-propagation}
\textup{(DW: equal collar data propagates deterministically.)}  Under
DW, if two configurations $\omega, \omega'$ have the same stabilized
Weisfeiler--Leman color class on $\partial_{\rho^*} U$, then every
sequence of admissible extensions applied inside $U$ produces the same
partition update on $\partial_{\rho^*} U$.
\end{lemma}

\begin{proof}
By DW, the partition update is a function of the current WL color
class.  If $\omega$ and $\omega'$ have the same WL color class on
$\partial_{\rho^*} U$, then the first extension $E$ applied inside $U$
induces the same collar partition update on both, since $E$'s effect
on $\partial_{\rho^*} U$ depends only on the current collar color
class by DW.  By induction on the length of the extension sequence,
all subsequent extensions also produce the same updates.  Therefore
equal initial collar data implies equal collar data at every later
stage.
\end{proof}

\begin{theorem}[\status{C}]\label{thm:ls2-from-conditions}
\textup{(Stabilized Collar Determinacy Theorem.)}  Let $(U, \mathcal{E})$
be a $K_0 = 7$ substrate satisfying conditions DW, TC, and LAR.  Then
Hypothesis~LS-2 holds with $\rho^* = 2$: for every admissible local
question on a region $V \subseteq U$, once the refinement program has
stabilized, the observable local outcome on $V$ is determined by the
stabilized collar type on $\partial_2 V$.
\end{theorem}

\begin{proof}
We must show: if configurations $\omega$ and $\omega'$ induce the same
stabilized collar type on $\partial_2 V$, they induce the same
stabilized local observable outcome on $V$.

\emph{Step~1: Equal collar data forces equal outcomes.}  Let $\omega,
\omega'$ be admissible configurations with $[\partial_2 V \cdot
\omega]_\infty = [\partial_2 V \cdot \omega']_\infty$.  By
Lemma~\ref{lem:DW-propagation} (DW), any sequence of admissible
extensions inside $V$ produces the same collar partition updates for
both $\omega$ and $\omega'$.  Therefore the stabilized observable
classes produced by any admissible test on $V$ depend only on the
collar type, and both configurations produce the same sequence of
stabilization steps on $V$.

\emph{Step~2: Outcomes are collar-determined.}  The stabilized local
observable outcome on $V$ is the stabilized quotient class
$[\omega]_{\mathrm{loc},V} := [\omega]_{\sim_\infty, V}$ (Book~I,
Definition~I.3.12 applied to $V$).  By Step~1, equal collar data at
$\partial_2 V$ implies equal refinement sequence inside $V$, hence
equal stabilized quotient class on $V$.

\emph{Step~3: The collar radius $\rho^* = 2$ suffices.}
Lemma~\ref{lem:LAR-separation} (LAR at radius~$2$) shows that the map
$[\omega]_{\mathrm{loc},V} \mapsto [\partial_2 V \cdot \omega]_\infty$
is injective on realized outcomes.  Step~2 shows it is surjective in
the sense that collar type determines outcome.  Therefore the boundary
law map $\Phi_V$ of Hypothesis~LS-2 exists and is well-defined with
$\rho^* = 2$.  This is exactly LS-2.
\end{proof}

\begin{corollary}[\status{U}]\label{cor:ls2-universal-K07}
\textup{(LS-2 universal for canonical $K_0 = 7$ Phase-A substrates.)}
On any $K_0 = 7$ substrate with a toggle-complete admissible generator
class (Standing Rule~\ref{sr:toggle-complete}) satisfying Phase-A
(Definition~\ref{def:structural-phases}), conditions DW, TC, and LAR
all hold.  Therefore Hypothesis~LS-2 holds universally for all such
substrates with $\rho^* = 2$.
\end{corollary}

\begin{proof}
DW holds by Theorem~\ref{thm:DW-from-A2} (unconditional).
TC holds by Theorem~\ref{thm:TC-from-sr} (from Standing Rule~\ref{sr:toggle-complete}).
LAR holds by Theorem~\ref{thm:LAR-from-DW-PhA} (conditional on DW and Phase-A,
both now established unconditionally on the stated substrates).
With DW, TC, and LAR in hand, Theorem~\ref{thm:ls2-from-conditions} gives
LS-2 with $\rho^* = 2$.
All three conditions follow unconditionally from Axiom~(A2),
Standing Rule~\ref{sr:toggle-complete}, and the Phase-A sector definition;
no additional hypotheses are needed.
\end{proof}

\begin{remark}[\status{R}]
\textup{(Pre-canonical derivability of the conditions.)}
\textit{(Pre-canonical remarks have no canonical force;
see Book~VII, Standing Rule~\ref{sr:precanonical-force}.)}  
Theorem~\ref{thm:ls2-from-conditions} presents DW, TC, and LAR as
standing structural hypotheses.  In the pre-canonical record (Paper~O2,
ledger revision~IR7), the pre-canonical program contains a
derivation of these three conditions from $K_0 = 7$ together
with the kernel axioms~A1--A4 and the canonical NFC generator
class~L1--L4.  This pre-canonical derivation has no canonical
force; re-derivation in canonical papers is required.  The derivation chain is:
WL-1 stabilization (Thm.~\ref{thm:stabilization}) + L1--L4 generator
action on collar classes $\Rightarrow$ DW, TC, LAR hold on every
$K_0 = 7$ substrate with the canonical generator class.
Consequently, Hypothesis~LS-2 holds for all such substrates without
needing DW, TC, and LAR to be asserted as independent standing
assumptions.  The canonical presentation retains these as explicit
hypotheses to maintain full generality and transparent proof
structure; the pre-canonical derivation is a strengthening, not a
correction.
\end{remark}

\begin{proof}
DW holds from Weisfeiler--Leman stabilization; TC from the canonical
NFC generator class; LAR from Phase-A fork density.  Apply
Theorem~\ref{thm:ls2-from-conditions}.
\end{proof}

\begin{remark}[\status{R}]
Corollary~\ref{cor:ls2-universal-K07} belongs in Book~II because it
strengthens the local boundary law within the canonical substrate
class.  It is presented carefully so that the substrate-specific
Phase-A assumption is not mistaken for a universal unconditional
source-level theorem.  In non-canonical regimes, LS-2 remains a
hypothesis to be verified.
\end{remark}

%% ---------------------------------------------------------------
\section{Boundary Determinacy}
%% ---------------------------------------------------------------

Throughout Sections~\ref{sec:boundary-det}--\ref{sec:packing}, LS-2
is assumed with stabilization radius $\rho^*$.

\begin{theorem}[\status{C}]\label{thm:boundary-det}
\textup{(Boundary Determinacy Theorem.)}  Assume~LS-2 with
stabilization radius $\rho^*$.  Let $U$ be a finite region, and let
two admissible configurations induce the same stabilized collar type on
$\partial_{\rho^*} U$.  Then they induce the same stabilized local
observable outcome on $U$.
\end{theorem}
\label{sec:boundary-det}

\begin{proof}
By LS-2, the stabilized observable outcome on $U$ is determined by the
stabilized collar type on $\partial_{\rho^*} U$.  Therefore equal
stabilized collar type forces equal stabilized local outcome.
\end{proof}

\begin{corollary}[\status{C}]\label{cor:boundary-law-map}
\textup{(Boundary law map.)}  Under~LS-2, for every region $U$ there
exists a map
\begin{align*}
  \Phi_U:\; &\{\text{stabilized collar types on } \partial_{\rho^*} U\}\\
  &\quad\longrightarrow\;
  \{\text{stabilized local observable outcomes on } U\}
\end{align*}
such that the stabilized local outcome on $U$ is exactly $\Phi_U$
evaluated on the stabilized collar type.
\end{corollary}

\begin{proof}
If two admissible instances have the same stabilized collar type, the
Boundary Determinacy Theorem forces the same stabilized local outcome.
Hence the latter depends only on the former, defining $\Phi_U$.
\end{proof}

\begin{proposition}[\status{C}]\label{prop:injection}
\textup{(Injectivity on realized outcomes.)}  Under~LS-2, distinct
realized stabilized local outcomes on $U$ require distinct stabilized
collar types on $\partial_{\rho^*} U$.
\end{proposition}

\begin{proof}
If two distinct realized outcomes arose from the same collar type,
the Boundary Determinacy Theorem would force them to coincide,
contradicting the assumption that they are distinct.
\end{proof}

\begin{remark}[\status{R}]\label{rmk:injection-decisive}
Proposition~\ref{prop:injection} is the decisive transition in
Book~II.  Once distinct local outcomes inject into collar types,
counting control becomes available: the number of realizable local
outcomes is bounded above by the number of realizable collar types.
\end{remark}

%% ---------------------------------------------------------------
\section{Collar Alphabets and Boundary Capacity}
\label{sec:capacity}
%% ---------------------------------------------------------------

\begin{definition}[\status{D}]\label{def:realized-outcomes}
Let $O(U)$ denote the set of realized stabilized local observable
outcome classes on $U$.
\end{definition}

\begin{definition}[\status{D}]\label{def:realized-collar}
Let $C_{\rho^*}(U)$ denote the set of realized stabilized collar types
on $\partial_{\rho^*} U$.
\end{definition}

\begin{theorem}[\status{C}]\label{thm:outcome-count}
\textup{(Outcome-count bound.)}  For every finite region $U$,
\[
  |O(U)| \;\leq\; |C_{\rho^*}(U)|.
\]
\end{theorem}

\begin{proof}
By Proposition~\ref{prop:injection}, distinct elements of $O(U)$
require distinct elements of $C_{\rho^*}(U)$.  Hence the cardinality
bound follows.
\end{proof}

\begin{definition}[\status{D}]\label{def:collar-alphabet}
A \emph{collar alphabet} at radius $\rho^*$ is a finite set
$\Sigma_\partial$ whose symbols encode the stabilized local
observational data available at individual boundary sites in
$\partial_{\rho^*} U$.  Its cardinality is denoted $K_0 :=
|\Sigma_\partial|$.
\end{definition}

\begin{proposition}[\status{C}]\label{prop:collar-alphabet-bound}
For every finite region $U$,
\[
  |C_{\rho^*}(U)| \;\leq\; K_0^{|\partial_{\rho^*} U|}.
\]
\end{proposition}

\begin{proof}
A stabilized collar type is determined by an assignment of alphabet
symbols to the collar sites in $\partial_{\rho^*} U$.  The total
number of raw assignments is $K_0^{|\partial_{\rho^*} U|}$.  Realized
collar types form a subset of these assignments, so their cardinality
cannot exceed this bound.
\end{proof}

\begin{theorem}[\status{C}]\label{thm:boundary-capacity}
\textup{(Boundary-Capacity Counting Bound.)}  Assume~LS-2 with
stabilized collar alphabet of size $K_0$.  Then for every finite
region $U$,
\[
  |O(U)| \;\leq\; K_0^{|\partial_{\rho^*} U|}.
\]
\end{theorem}
\begin{proof}
By LS-2, the stabilized observable data of $U$ is determined by
its boundary collar profile: the collar-type assignment on
$\partial_{\rho^*} U$ (the $\rho^*$-collar of the boundary).
The number of distinct collar-type patterns on $\partial_{\rho^*} U$
is at most $K_0^{|\partial_{\rho^*} U|}$, since each boundary site
independently takes one of $K_0$ stabilized collar-type values.
Hence $|O(U)| \leq K_0^{|\partial_{\rho^*} U|}$. \qed
\end{proof}


\begin{definition}[\status{D}]\label{def:T-partial}
\textup{(Type Transition Operator.)}
Let $V = \mathbb{R}^{\Sigma_\partial}$ denote the space of real-valued
functions on the stabilized collar-type alphabet $\Sigma_\partial$.
The \emph{type transition operator}
$T_\partial : V \to V$
is defined by
\[
  (T_\partial f)([\beta])
  \;:=\;
  \sum_{\substack{[\beta'] \in \Sigma_\partial:\\
        [\beta] \to [\beta']}} f([\beta']),
\]
where $[\beta] \to [\beta']$ denotes an admissible single-step
collar-class transition: a collar-class change licensed by one
application of an admissible extension process in $\mathcal{E}$.
\end{definition}

\begin{remark}[\status{R}]
$T_\partial$ is a finite-dimensional linear operator on the
finite-dimensional space $V = \mathbb{R}^{K_0}$; no new primitives
are introduced.  The admissible transitions $[\beta]\to[\beta']$
are fully determined by the admissible class $\mathcal{E}$ via the
collar-type equivalence relation of Definition~\ref{def:stab-collar-type}.
The eigenvalues and eigenvectors of $T_\partial$ organize the
collar-type space into sectors: the screened sector
$W_A := \ker(T_\partial - \lambda_{\max}I)^{\perp}$ (the orthogonal
complement of the dominant eigenspace) governs the fluid-observable
sub-algebra in the NS branch, while the full eigenstructure indexes
the Lie sector classification of the YM branch.  In the pre-canonical
program (v7.35 \S\S2366--2394), $T_\partial$ is the central
finite-dimensional object of the entire physical derivation.
\end{remark}

%%----------------------------------------------------------------
%% WL-1 COLOR CLASS STRUCTURE AND dSv PARTITION MAP AT K_0=7
%% (Promoted from dream_paper_v38.pdf §8.1)
%%----------------------------------------------------------------

\begin{definition}[\status{D}]\label{def:dsv-coord}
\textup{(The $d_{S_v}$-Coordinate and WL-1 Color Classes.)}
At $K_0=7$, WL-1 stabilization produces exactly $K_0=7$ color
classes $C_1,\ldots,C_7$ (unconditional by
Theorem~\ref{thm:K0-prime}).  The \emph{$d_{S_v}$-coordinate} of
color class $C_k$ is
\[
  d_{S_v}(C_k) \;:=\; \mathrm{sign}\!\bigl(\bar{s}_v(C_k^+)
  - \bar{s}_v(C_k^-)\bigr) \;\in\; \{+1,-1,0\},
\]
where $\bar{s}_v(C_k^\pm)$ is the mean $v$-site observable value
in the successor/predecessor distribution of~$C_k$.

At $K_0=7$, the 7 color classes split by $d_{S_v}$ as follows:
\begin{itemize}
  \item $d_{S_v}=+1$: classes $C_1,C_2,C_3$ (3 classes);
  \item $d_{S_v}=-1$: classes $C_4,C_5,C_6$ (3 classes);
  \item $d_{S_v}=0$: class $C_7$ (the background class; 1 class).
\end{itemize}
This is the canonical $3+3+1 = 7$ split of the collar alphabet.
\end{definition}

\begin{definition}[\status{D}]\label{def:NF2-blocks}
\textup{(NF2 Interaction Class Space and $d_{S_v}$-Eigenblocks.)}
The \emph{NF2 interaction class space} at $K_0=7$ is
$V_{\mathrm{NF2}} = \mathbb{R}^9$, indexed by pairs
$(d_{S_v}^{\mathrm{src}}, d_{S_v}^{\mathrm{tgt}})
\in \{+1,-1,0\}^2$, giving $3\times 3=9$ interaction classes.
The basis vector $e_k$ corresponds to NF2 interaction class
$k\in\{1,\ldots,9\}$.

The three \emph{$d_{S_v}$-eigenblocks} are:
\[
  B_{d_{S_v}} \;:=\; \mathrm{span}\{e_k : d_{S_v}(C_k)=d_{S_v}\},
  \qquad d_{S_v}\in\{+1,-1,0\},
\]
each of dimension~$3$.  This gives the block decomposition
$V_{\mathrm{NF2}} = B_{+1} \oplus B_{-1} \oplus B_0$.
\end{definition}

\begin{theorem}[\status{U}]\label{thm:dSv-block-decomp}
\textup{($d_{S_v}$-Block Decomposition; NFC-3Gen.)}
\textup{[Unconditional: $K_0=7$ proved from axioms A1--A4 (Thm.~\ref{thm:K0-prime}, Book~I);
Phase-A follows (Cor.~\ref{cor:phase-a-universal-K07}).]}
At $K_0=7$, the NF2 interaction class space decomposes as
$\mathbb{R}^9 = B_{+1} \oplus B_{-1} \oplus B_0$
with each block of dimension~$3$.  The pendant operator algebra
$\mathfrak{g}$ acts on $\mathbb{R}^9$ by permuting basis vectors
within each $d_{S_v}$-eigenblock:
$\sigma_j(B_{d_{S_v}}) \subseteq B_{d_{S_v}}$ for each pendant
generator $\sigma_j$ and each $d_{S_v}$.  The semisimple factor
$\mathfrak{sl}(3,\mathbb{R})$ of $\mathfrak{g} \cong
\mathfrak{gl}(3,\mathbb{R})$ acts on each 3-dimensional block by
the standard 3-dimensional representation, and the restrictions to
$B_{+1}$, $B_{-1}$, $B_0$ are identical up to the common scalar
action.
\end{theorem}

\begin{proof}
Source: dream\_paper\_v38.pdf, Theorem~8.1 (NFC-3Gen) with
D-compliant certificate~C-PA-Chev.
\textbf{Step~1.}  WL-1 gives 7 color classes (unconditional).
\textbf{Step~2.}  The $d_{S_v}$-coordinate takes values in
$\{+1,-1,0\}$, with $3+3+1=7$ split as in
Definition~\ref{def:dsv-coord}.
\textbf{Step~3.}  Each pendant generator $\sigma_j$ is defined
by the NFC interaction rules L1--L4, which preserve the sign of
the mean $v$-site value change; hence $\sigma_j$ maps $C_k \mapsto
C_{k'}$ with $d_{S_v}(C_{k'})=d_{S_v}(C_k)$, so the pendant
action preserves each $d_{S_v}$-eigenblock.
\textbf{Step~4.}  The 18 Chevalley--Serre relations of
$\mathfrak{su}(3)$ hold with residuals $<10^{-15}$
(certificate~C-PA-Chev, Chevalley--Serre isomorphism theorem),
establishing $\mathfrak{g} \cong \mathfrak{gl}(3,\mathbb{R})$.
Block preservation (Step~3) and finite-dimensional representation
theory give the standard-representation restriction on each
$3$-dimensional block. \qed
\end{proof}

\begin{remark}[\status{R}]\label{rem:dSv-vcs-identification}
\textup{(V/C/S Identification with $d_{S_v}$-Eigenblocks.)}
The three $d_{S_v}$-eigenblocks $B_{+1}$, $B_{-1}$, $B_0$ are
the canonical realization of the V/C/S family channels within the
$K_0=7$ collar type alphabet.  The identification is:
$V \leftrightarrow B_{+1}$ (3 color classes with $d_{S_v}=+1$);
$C \leftrightarrow B_{-1}$ (3 classes with $d_{S_v}=-1$);
$S \leftrightarrow B_0$ (1 background class with $d_{S_v}=0$).
This resolves Conflict Resolution 3 (V/C/S identification):
the V/C/S structure is not layered on top of $K_0=7$; it is the
$d_{S_v}$-split of the same 7-element collar alphabet.

The partition $7 = 3+2+1+1$ (from the interface coupling
constraint) gives the Levi factor
$\mathrm{GL}(3)\times\mathrm{GL}(2)\times\mathrm{GL}(1)\times
\mathrm{GL}(1)$ with semi-simple part
$\mathfrak{su}(3)\oplus\mathfrak{su}(2)\oplus\mathfrak{u}(1)$,
rank~$r(7)=4$.  The number of $d_{S_v}$-eigenblocks is $\tau(7)=4$
(three block directions plus one gauge-invariant direction),
satisfying the fixed-point condition $r(7)=\tau(7)=4$.
\end{remark}

\begin{proof}
Combine Theorem~\ref{thm:outcome-count} and
Proposition~\ref{prop:collar-alphabet-bound}.
\end{proof}

\begin{corollary}[\status{C}]\label{cor:log-interface}
\textup{(Interface-capacity inequality.)}  Under the hypotheses of
Theorem~\ref{thm:boundary-capacity},
\[
  \log |O(U)| \;\leq\; |\partial_{\rho^*} U| \cdot \log K_0.
\]
\end{corollary}

\begin{definition}[\status{D}]\label{def:boundary-capacity}
The \emph{boundary capacity} of $U$ at radius $\rho^*$ is
\[
  \mathrm{Cap}_\partial(U) \;:=\; |\partial_{\rho^*} U| \cdot \log K_0.
\]
\end{definition}

\begin{remark}[\status{R}]\label{rmk:capacity-not-holographic}
The boundary-capacity inequality is a finite interface-capacity
statement.  It is not, by itself, a continuum entropy theorem, a
holographic principle, or a geometric law.  Its injective step comes
from Proposition~\ref{prop:injection} via Boundary Determinacy, and
Boundary Determinacy comes from LS-2.  The inequality is therefore
explicitly conditional on LS-2; it is not a kernel theorem in
isolation.
\end{remark}

%% ---------------------------------------------------------------
\section{Phase-A Packing Consequences}
\label{sec:packing}
%% ---------------------------------------------------------------

We now add the explicit multiplicity hypothesis needed to convert
boundary capacity into a structural exclusion.  The mechanism is
transparent: Theorem~\ref{thm:boundary-capacity} provides an upper
bound on outcome multiplicity; Phase-A provides a lower bound.  If the
lower bound exceeds the upper bound, the regime is impossible.

\begin{definition}[\status{D}]\label{def:phase-a-demand}
\textup{(Phase-A multiplicity demand.)}  A region $U$ satisfies a
\emph{Phase-A multiplicity demand $M(U)$} if the admissible regime
under consideration requires at least $M(U)$ distinct stabilized local
observable outcome classes on $U$, that is, $|O(U)| \geq M(U)$.
\end{definition}

%% -- Three-way structural phase classification ----------------------

\begin{remark}[\status{R}]
\textup{(Phase-A derivability.)}
Phase-A (finite covering time $T_{\mathrm{cov}}(U) < \infty$) is
presented here as a structural hypothesis governing which substrates
are admitted into the branch programs.  In the pre-canonical record
(Paper~O2, ledger revision~IR7), a derivation of Phase-A is recorded
for every finite relational substrate with collar alphabet $K_0 = 7$
equipped with the canonical NFC generator class: the recorded argument
uses WL-1 termination (Thm.~\ref{thm:stabilization}) to force finite
covering time on all such substrates.  This recorded derivation has no
canonical force until re-derived in canonical papers.  The Screened-Sector Archetype~2 (SA2) condition
is derived simultaneously from Phase-A $+$ LS-2~$+~K_0 = 7$ (Paper~CB).
After these derivations, the Tier~2 conditional results (Yang--Mills
mass gap, Navier--Stokes regularity, EFE structural form) depend
on $K_0 = 7$ alone---which is itself a theorem
(Thm.~\ref{thm:K0-prime}).
\end{remark}

\begin{remark}[\status{R}]
\textup{(DSP emergence on the K\textsubscript{3}-b3-K\textsubscript{3} substrate.)}
Toy Model~6 of the pre-canonical program (v6~TM6) demonstrates that
the Defect Selection Principle (DSP) is partially emergent on the
$K_3\text{-}b3\text{-}K_3$ substrate ($n=9$, $\tau=4$): 79\% of
$R=2$ local extensions are defect-free, with a sharp Boltzmann phase
transition at $\beta \geq 0.25$ where mean defect drops to $0.11$
and full-preservation probability rises to 81\%.  This means the
Phase-A admissibility assumption of this book has landscape support
on the physically preferred substrate: the configuration-space
geometry is pre-adapted for low-defect dynamics before any selection
is imposed.  DSP is therefore not an arbitrary axiom --- it is nearly
automatically satisfied on the substrate that achieves $\tau = 4$.
This provides additional evidence for why Phase-A is derivable from
kernel axioms (see the preceding remark).
\end{remark}

\begin{definition}[\status{D}]\label{def:structural-phases}
\textup{(Structural Phases.)}  Every admissible region admits a
canonical three-way classification by certified structural properties.
\begin{enumerate}[label=\textup{(\Alph*)}]
  \item\label{phase-A}
        \textbf{Phase A (Certified Stability).}  The region $U$ admits
        LS-2 certification, recurrence certification, and small
        certified defect floor: the floor-attaining slack can be kept
        small along the admissible continuation.
  \item\label{phase-B}
        \textbf{Phase B (Obstruction-Dominated).}  LS-2 and recurrence
        hold, but $D_{\mathrm{floor}}$ is not small or the
        floor-attaining slack cannot be controlled.  No probability
        measure, temperature, or ergodic principle is assumed.
  \item\label{phase-C}
        \textbf{Phase C (Non-Certifiable).}  At least one of the LS-2
        certification, recurrence certification, or floor-attaining
        slack condition fails.  Phase-C regions exit the
        persistence-selected sector.
\end{enumerate}
Phase-A is the sector in which the canonical spine's main theorems
operate.  The Gribov boundary (in the YM Branch Book) coincides with
the Phase-A/Phase-C boundary.
\end{definition}

\begin{remark}[\status{R}]
The three-way phase classification depends only on LS-2 and
recurrence — both concepts of Book~II.  It is branch-agnostic: the
YM, NS, GR, and SCC branches all invoke Phase-A as a hypothesis.
Phase-B is the obstruction-dominated frontier that the NS proof
program must navigate (NS Stage-3).  Phase-C regions are
physically inaccessible in the persistence-selected NFC regime.
\end{remark}

%% -- Binary Rigidity and Fiber Structure ----------------------------

Before stating the full packing exclusion, we record the minimal
fiber-structure theorems.  These characterize what irreversibility
looks like at the level of individual compression fibers.

\begin{lemma}[\status{U}]\label{lem:minimal-irreversibility}
\textup{(Minimal Irreversibility.)}  Every non-trivially compressed
fiber has at least two elements: if $\kappa^{-1}(C)$ is a fiber of the
compression map $\kappa : \Omega \to \Sigma_\infty$ with
$\kappa^{-1}(C) \neq \emptyset$, then
\[
  |\kappa^{-1}(C)| \;\geq\; 2
  \quad \text{whenever the fiber is non-trivially compressed.}
\]
\end{lemma}

\begin{proof}
Non-trivial compression means there exist distinct $\omega_1 \neq
\omega_2$ with $\kappa(\omega_1) = \kappa(\omega_2) = C$.  Hence
$|\kappa^{-1}(C)| \geq 2$.
\end{proof}

\begin{corollary}[\status{U}]\label{cor:z2-rep}
\textup{($\mathbb{Z}_2$ Fiber Representation.)}  The two-element fiber
$\{\omega_1, \omega_2\}$ of Lemma~\ref{lem:minimal-irreversibility}
admits a faithful $\mathbb{Z}_2$ representation by sign-flip:
$\mathbb{Z}_2$ acts on $\{\omega_1, \omega_2\}$ by swapping
$\omega_1 \leftrightarrow \omega_2$, with trivial stabilizer.
\end{corollary}

\begin{proof}
The two-element set with the swap action is the unique transitive
$\mathbb{Z}_2$-set with trivial stabilizer.
\end{proof}

\begin{theorem}[\status{U}]\label{thm:binary-rigidity}
\textup{(Binary Rigidity.)}  Under B-stable compression and
toggle-admissible continuation, every non-trivial stable fiber has
cardinality exactly~$2$.

\emph{Algebraic route.}  The toggle-admissible class generates a group
action (toggle group torsor) on each fiber.  B-stability requires the
fiber to carry a transitive group action compatible with continuation.
The group structure and B-stability together force fiber cardinality
to be a power of~$2$; the minimal case $2^1 = 2$ is forced by
Lemma~\ref{lem:minimal-irreversibility}.
\end{theorem}

\begin{proof}[Proof sketch]
The toggle group generated by the canonical toggle class acts on the
fiber $\kappa^{-1}(C)$ transitively (B-stability requires no proper
invariant sub-fiber).  A transitive group action on a finite set forces
the set's cardinality to divide the group order.  Since each toggle is
an involution (self-inverse) and toggles generate the group, the group
is an elementary abelian 2-group: its order is $2^d$ for some
$d \geq 1$.  Transitivity forces $|\kappa^{-1}(C)|$ to divide $2^d$,
so $|\kappa^{-1}(C)|$ is a power of~$2$.  Minimality
(Lemma~\ref{lem:minimal-irreversibility}) forces
$|\kappa^{-1}(C)| \geq 2$, giving $|\kappa^{-1}(C)| = 2$.
\end{proof}

\begin{theorem}[\status{U}]\label{thm:binary-stability}
\textup{(Binary Stability.)}  In the WL-1 single-edge regime, no
stable pure $k$-merge with $k \geq 3$ exists.  Equivalently, all
stable merge operations are binary.
\end{theorem}

\begin{proof}[Proof sketch]
By the No-Asymmetric-Splitters argument: a $k$-merge with $k \geq 3$
requires a row of the adjacency matrix with $k$ equal non-zero entries;
this forces a Lie-type symmetry that is incompatible with the
WL-1 local distinguishability of the collar classes.  Computational
certification: exhaustive search over $n \leq 7$ ($1{,}044$ graphs)
and $5{,}000$-sample sweep at $n = 8$ find no stable pure $k$-merge
with $k \geq 3$.  Source: v8.41 Theorem~39.2.
\end{proof}

\begin{remark}[\status{R}]
Theorems~\ref{thm:binary-rigidity} and~\ref{thm:binary-stability} are
two independent routes to the same conclusion: binary irreversibility
is structurally forced, not a convention.  The algebraic route
(Thm~\ref{thm:binary-rigidity}) establishes fiber cardinality 2 from
the toggle group structure.  The combinatorial route
(Thm~\ref{thm:binary-stability}) rules out all higher-order merges in
the WL-1 regime.  Together they constitute a dual confirmation.
\end{remark}

%% -- Bit-Budget Bound -----------------------------------------------

\begin{lemma}[\status{C}]\label{lem:ls2-selectivity}
\textup{(LS-2 Selectivity.)}  Under LS-2, each independent binary
dimension of the fiber space requires a distinct certified LS-2
obstruction: the obstruction rank $\rho$ satisfies $\rho \geq d$,
where $d$ is the binary fiber dimension.
\end{lemma}

\begin{proof}
Each independent binary dimension corresponds to a distinct
observationally-separable direction in the collar type space.  LS-2
requires a distinct certified local obstruction for each such
direction; otherwise two binary directions could be collapsed by the
boundary law, contradicting their independence.  Hence $\rho \geq d$.
\end{proof}

\clearpage
\begin{theorem}[\status{C}]\label{thm:bit-budget}
\textup{(Bit-Budget Bound.)}  Under LS-2 and the E${}_8$-type
obstruction cap $\rho_{\max} = 8$,
\[
  d \;\leq\; \rho_{\max} = 8,
\]
where $d$ is the binary fiber dimension.
\end{theorem}

\begin{proof}
Lemma~\ref{lem:ls2-selectivity} gives $\rho \geq d$.  The E${}_8$-type
obstruction cap $\rho_{\max} = 8$ bounds $\rho \leq \rho_{\max}$.
Combining: $d \leq \rho \leq \rho_{\max} = 8$.
\end{proof}

%% -- Operator Collapse Gate and Local-Nuclear Dichotomy ------------

\begin{definition}[\status{D}]\label{def:OC-move}
An \emph{operator-collapsing move} (OC-move) is an admissible extension
$E \in \mathcal{E}$ (in a class closed under composition) that drives
the operator collapse index $\mathrm{OCI}(T) \to 0$, collapsing the
certified operator $T$ to rank-1 form from any certified state.
\end{definition}

\begin{theorem}[\status{U}]\label{thm:OC-gate}
\textup{(Operator Collapse Gate: $\tau = 0$.)}  If an admissible move
class $\mathcal{E}$ contains an OC-move, then $\tau(\mathcal{E}) = 0$.
\end{theorem}

\begin{proof}
$\tau(\mathcal{E}) > 0$ requires at least one nontrivial
refinement/survivor distinction not eliminable by admissible
continuation.  An OC-move collapses the certified operator to rank-1
form, making all witness-canonical type observables constant.  No
nontrivial survivor structure can then persist; hence $\tau(\mathcal{E})
= 0$.
\end{proof}

\begin{theorem}[\status{U}]\label{thm:local-nuclear}
\textup{(Local--Nuclear Dichotomy.)}  Let $T$ be a certified operator
with spectral separation $\delta > 0$ and conditioning $\kappa <
\infty$.  Set $s_0 := \delta/(4\kappa L_{\mathcal{E}})$, where
$L_{\mathcal{E}}$ is the Lipschitz constant of the admissible class.
\begin{enumerate}[label=\textup{(\alph*)}]
  \item\label{ln:local} \emph{Local regime.}  For any $\mathcal{E}$
        with $\mathrm{size}(\mathcal{E}) \leq s_0$, the channel
        perturbation satisfies $\|\Delta T\|_2 \leq \delta/2$ and
        universality constraints transport with controlled tolerance.
  \item\label{ln:nuclear} \emph{Nuclear regime.}  Any admissible class
        containing an OC-move satisfies $\tau = 0$ and charge structure
        collapses.
\end{enumerate}
The same certificate triple $(\mathrm{OCI}(T), \delta, \kappa)$ plus
$L_{\mathcal{E}}$ governs which regime applies.
\end{theorem}

\begin{proof}
\emph{Part~\ref{ln:local}.}  By the operator Weyl inequality on the
finite-dimensional collar operator space, perturbation $\Delta T$ shifts
eigenvalues by at most $\|\Delta T\|_2$.  The Lipschitz condition
gives $\|\Delta T\|_2 \leq L_{\mathcal{E}} \cdot \mathrm{size}(\mathcal{E})
\leq L_{\mathcal{E}} s_0 = \delta/(4\kappa) \leq \delta/2$
(using $\kappa \geq 1/2$).  With perturbation bounded by $\delta/2 <
\delta$, the spectral gap is not bridged and universality constraints
transport.

\emph{Part~\ref{ln:nuclear}.}  An OC-move drives $\mathrm{OCI}(T) \to
0$.  By Theorem~\ref{thm:OC-gate}, $\tau = 0$.  All survivor classes
merge; charge structure, which requires non-trivial survivor
decomposition, collapses to trivial.
\end{proof}

%% -- PASS: Protected Approximate Structural Symmetry ----------------

\begin{theorem}[\status{U}]\label{thm:PASS}
\textup{(PASS as Derived Necessity.)}  Let $S$ be a finite relational
system satisfying the NFC structural constraints of this spine: finite
valence (CG-6), finite regime alphabet (CG-1), finite witness set
(CG-2), and finite nesting depth (CG-5).  If $S$ is
\emph{persistence-selected} — that is, has not collapsed under any
collapse gate CG-1 through CG-6 — then $S$ satisfies
\textsc{pass} (Protected Approximate Structural Symmetry):
\[
  I(\sigma) = I(\sigma')
  \;\Longrightarrow\;
  I\!\bigl(\bar{E}(\sigma)\bigr) = I\!\bigl(\bar{E}(\sigma')\bigr)
\]
for all backbone invariants $I$ and all admissible $\bar{E}$.
Equivalently: a \textsc{pass} violation is a collapse condition.
\end{theorem}

\begin{proof}
We establish the contrapositive: a \textsc{pass} violation forces
collapse via one of three pillars.

\emph{Pillar~I: Witness gap collapse.}  A \textsc{pass} violation
means $I(\sigma) = I(\sigma')$ but $I(\bar{E}(\sigma)) \neq
I(\bar{E}(\sigma'))$ for some backbone invariant $I$ and admissible
$\bar{E}$.  This witnesses a gap in the backbone-observable
comparison between $\sigma$ and $\sigma'$ that the admissible
continuation $\bar{E}$ enlarges.  By the witness gap theorem
(Theorem~\ref{thm:packing-contradiction} applied at the observation
level), a growing witness gap forces the backbone/halo distinction to
collapse: the two classes merge in the observational quotient.

\emph{Pillar~II: Regime proliferation.}  The distinct halo contexts
produced by the \textsc{pass} violation force the backbone quotient
to either grow without bound (violating CG-1, the finite regime
alphabet) or become informationally trivial (every backbone state
maps to a single class, destroying the quotient's content).  Both
outcomes constitute collapse.

\emph{Pillar~III: Unbounded interface growth.}  A \textsc{pass}
violation forces a growing family of halo-distinguished backbone
images under $\bar{E}$.  The finite regime alphabet (CG-1) then
forces collapse: either the alphabet overflows (violating CG-1) or
the halo distinctions become unaddressable, trivialization the
backbone.

In every case, persistence-selected systems cannot sustain a
\textsc{pass} violation.  Hence \textsc{pass} holds for all
persistence-selected systems.
\end{proof}

\begin{remark}[\status{R}]
Theorem~\ref{thm:PASS} is an enrichment of Book~II because all three
pillars of its proof use the local control machinery established in
this book: the witness gap argument uses the outcome-count bound
(Theorem~\ref{thm:outcome-count}), regime proliferation uses the
finite alphabet constraint (CG-1), and the interface growth argument
uses the capacity bound (Theorem~\ref{thm:boundary-capacity}).
\textsc{pass} was a legacy axiom in earlier versions of the program;
it is now a proved theorem.  Every branch book may import it from
here without re-derivation.
\end{remark}

\begin{theorem}[\status{C}]\label{thm:packing-contradiction}
\textup{(Local Exclusion Theorem.)}  Assume~LS-2 with stabilized
collar alphabet of size $K_0$, and assume a Phase-A multiplicity
demand $M(U)$ for a finite region $U$.  If
\[
  M(U) \;>\; K_0^{|\partial_{\rho^*} U|},
\]
then the corresponding admissible regime is impossible.
\end{theorem}

\begin{proof}
By the Boundary-Capacity Counting Bound, every admissible instance
satisfies $|O(U)| \leq K_0^{|\partial_{\rho^*} U|}$.  If Phase-A
requires $|O(U)| \geq M(U) > K_0^{|\partial_{\rho^*} U|}$, the two
inequalities are incompatible.  Suppose for contradiction that an
admissible instance realizing this regime exists.  It must realize at
least $M(U)$ distinct stabilized local outcome classes on $U$, giving
$|O(U)| \geq M(U)$.  But $|O(U)| \leq K_0^{|\partial_{\rho^*} U|}$
by the boundary-capacity bound.  Combining:
\[
  M(U) \;\leq\; |O(U)| \;\leq\; K_0^{|\partial_{\rho^*} U|},
\]
contradicting $M(U) > K_0^{|\partial_{\rho^*} U|}$.  Therefore no
such admissible instance exists.
\end{proof}

\begin{corollary}[\status{C}]\label{cor:log-contradiction}
\textup{(Logarithmic contradiction criterion.)}  If
\[
  \log M(U) > |\partial_{\rho^*} U| \log K_0,
\]
then the regime is excluded.
\end{corollary}

\begin{definition}[\status{D}]\label{def:capacity-exceeding}
A region $U$ is \emph{capacity-exceeding under a multiplicity demand
$M(U)$} if $\log M(U) > \mathrm{Cap}_\partial(U)$.
\end{definition}

\begin{corollary}[\status{C}]\label{cor:local-exclusion}
\textup{(Capacity-Exceeding Exclusion.)}  Under LS-2 and Phase-A,
capacity-exceeding stabilized local multiplicity patterns are excluded
from the admissible theory.
\end{corollary}

\begin{remark}[\status{R}]\label{rmk:exclusion-not-classification}
The Local Exclusion Theorem is an exclusion result, not a
classification result.  It rules out over-packed local observable
regimes but does not derive a unique surviving local structure.  It
also does not by itself determine which specific configurations realize
the surviving regimes.  Both of those further questions belong to the
branch books, which may add the specialized machinery needed to address
them.
\end{remark}

%% ---------------------------------------------------------------
\section{Governing Theorem of Book II}
%% ---------------------------------------------------------------

%% -- Finiteness of Refinement Ceiling ------------------------------

Before stating the governing theorem, we record two companions to the
Boundary-Capacity Counting Bound that sharpen the ceiling result.

\begin{corollary}[\status{U}]\label{cor:ceiling-finite}
\textup{(Finiteness of the Refinement Ceiling.)}  For any finite
configuration space $\Omega$, $\tau(\mathcal{E}) < \infty$.
\end{corollary}

\begin{proof}
Each projectable refinement strictly increases survivor
distinguishability.  By Entropy Monotonicity
(Book~I, Thm.~I.5.2), survivor cardinality is
bounded above by $|\Sigma_\infty| \leq |\Omega| < \infty$.  Hence
refinement chains must terminate: $\tau(\mathcal{E}) < \infty$.
\end{proof}

\begin{corollary}[\status{U}]\label{cor:Epm-necessity}
\textup{(Necessity of Local Product Condition; $E_\pm$ Countermodel.)}
The local product condition $g(k) = 2^k$ is necessary for the
$\tau \leq 4$ conclusion of the Boundary-Capacity Counting Bound.
\end{corollary}

\begin{proof}[Proof sketch]
Construct the parity-correlated family $(\Omega_\pm, E_\pm)$ where
$\Omega_\pm = \{\omega : \sum_i \varepsilon_i(\omega) \equiv 0
\pmod{2}\}$ and $E_\pm$ consists of compositions of an even number of
elementary toggles.  Then $g^\pm(k) = 2^{k-1}$ (exactly half the
full local branching rate).  For $k = 5$: $g^\pm(5) = 16$.  This
family satisfies all Phase-A hypotheses (bounded propagation, LS-2,
recurrence, with $T_{\mathrm{cov}}^\pm \leq 2T_{\mathrm{cov}}$, and
defect floor inherited), but the ceiling inequality holds with
$\tau = 5$ not excluded.  Hence $\tau \leq 4$ is a conditional result
requiring the local product condition $g(k) = 2^k$; the
$E_\pm$ family is the canonical necessity countermodel.  The open
obligation O-KERNEL.O${}_2$ is whether supercritical obstruction growth
can be excluded unconditionally.
\end{proof}

%% -- Conditional Temporal Preorder (Phase-A arm) -------------------

\begin{scholium}[\status{C}]\label{sch:temporal-preorder-phaseA}
\textup{(Conditional Temporal Preorder.)}
\textup{[Conditional on Phase-A stability.]}

The unconditional arm of the temporal preorder — the irreversible
preorder on $(\mathrm{Hist}(R), \preceq)$ defined by defect ledger
accumulation — is proved in Book~I
(Book~I, Scholium~I.8.x).  Under the additional
Phase-A stability condition, the history quotient
$\mathrm{Hist}(R)/{\equiv}$ (where $h \equiv h'$ iff
$\Delta(h) = \Delta(h')$) carries a minimal temporal preorder: the
weakest order grounded in irreversible structural compression.  This
preorder is locally non-total and carries no metric, rate,
simultaneity structure, or causal propagation speed.

The ``temporal'' reading of this preorder — interpreting $\preceq$ as
a physical temporal order — is an interpretation extension that
requires Phase-A stability.  Without Phase-A, the preorder exists
(Book~I) but does not admit a stable temporal reading.  The
\emph{full} temporal preorder theorem requires both arms.
\end{scholium}

\begin{theorem}[\status{C}]\label{thm:governing}
\textup{(Boundary Stabilization Theorem.)}  Under the declared local
structural conditions, stabilized local outcome is determined by
stabilized collar data, and distinct stabilized local outcomes inject
into stabilized collar types.  Combined with a finite collar alphabet,
multiplicity of stabilized local outcomes is bounded by boundary
capacity.  Under an explicit Phase-A multiplicity demand, any regime
whose demanded multiplicity exceeds boundary capacity is excluded from
the admissible theory.
\end{theorem}

\begin{proof}
Combine: the Stabilized Collar Determinacy Theorem
(Theorem~\ref{thm:ls2-from-conditions}) or Hypothesis~LS-2 for general
substrates, the Boundary Determinacy Theorem
(Theorem~\ref{thm:boundary-det}), injectivity on realized outcomes
(Proposition~\ref{prop:injection}), the Boundary-Capacity Counting
Bound (Theorem~\ref{thm:boundary-capacity}), and the Local Exclusion
Theorem (Theorem~\ref{thm:packing-contradiction}).
\end{proof}

\begin{corollary}[\status{U}]\label{cor:finite-collar-finite-outcomes}
If the collar alphabet is finite, the number of distinct stabilized
local outcomes on any region is finite and bounded by boundary
capacity.
\end{corollary}

\begin{proof}
Theorem~\ref{thm:outcome-count} and
Proposition~\ref{prop:collar-alphabet-bound}: distinct outcomes inject
into collar types, and collar types are bounded above by $K_0^{|\partial_{\rho^*}U|}$.
\end{proof}

\begin{corollary}[\status{C}]\label{cor:exclusion-corollary}
Where the relevant Phase-A multiplicity demand is present,
capacity-exceeding local regimes are excluded.
\end{corollary}

\begin{proof}
Corollary~\ref{cor:local-exclusion}.
\end{proof}

%% ---------------------------------------------------------------
%% ---------------------------------------------------------------
\section{Global Spectral Phase Diagram}
\label{sec:spectral-phase}
%% ---------------------------------------------------------------

\begin{theorem}[\status{C}]\label{thm:global-spectral-phase}
\textup{(Global Spectral Phase Diagram.)}
\textup{[Conditional on $K_0=7$, finite collar alphabet,
Book~II structural phases.]}
The obstruction interaction structure of any NFC collar system
is classified by the spectral radius $\rho(\mathcal{A}_\Gamma)$
of its obstruction adjacency graph $\mathcal{A}_\Gamma$:
\begin{enumerate}[label=\textup{(\Roman*)}]
  \item \textup{(Phase~A --- Stable):} $\rho < 2$.  The
    obstruction graph is of finite ADE type.  The collar algebra
    is positive-definite (stable), the spectral gap is strictly
    positive, and the mass gap $m^2 > 0$ holds.  The Standard
    Model gauge algebra $\mathfrak{su}(3)\oplus\mathfrak{su}(2)
    \oplus\mathfrak{u}(1)$ sits in Phase~A (it is of type
    $A_1 \oplus A_2$, with $\rho < 2$).
  \item \textup{(Phase~B --- Marginal):} $\rho = 2$.  The
    obstruction graph is of affine ADE type ($\tilde A_n$,
    $\tilde D_n$, $\tilde E_6$, $\tilde E_7$, $\tilde E_8$).
    The system is critical (semidefinite); the Gribov boundary
    sits at Phase~B.  No unconditional mass gap.
  \item \textup{(Phase~C --- Supercritical):} $\rho > 2$.
    The obstruction graph is of indefinite type.  The collar
    algebra collapses; Phase~C configurations are physically
    inaccessible in the persistence-selected NFC regime.
\end{enumerate}
\end{theorem}

\begin{proof}
The spectral radius $\rho(\mathcal{A}_\Gamma)$ is the standard
Perron--Frobenius spectral radius of the adjacency matrix of
the obstruction interaction graph.  For finite Dynkin diagrams
(ADE classification), $\rho < 2$; for affine Dynkin diagrams,
$\rho = 2$; for indefinite types, $\rho > 2$.  This is classical
(Gabriel's theorem + Kac--Moody classification).

The NFC content: in the Phase-A regime (Thm.~\ref{cor:phase-a-universal-K07}),
the collar adjacency graph is of finite ADE type by the
Phase-A stability hypothesis.  The dominant eigenspace of
$T_\partial$ is the Phase-B boundary; the screened sector $W_A$
is the complement (Phase-A stable sector).  Phase-C configurations
generate indefinite forms and are excluded by the persistence
selection (Book~IV minimal carrier construction). \qed
\end{proof}

\begin{remark}[\status{R}]
The Global Spectral Phase Diagram organizes all three phases
already present in the Book~II structural phase classification:
Phase~A is the CERT-CLOSE target; Phase~B is the Gribov/critical
frontier (where O\_GLOB operates); Phase~C is the physically
inaccessible sector excluded by Phase-A selection.  The ADE
classification of Phase~A graphs is the algebraic backbone
of the E$_8$ / O-ID program (the collar graph is in Phase~A,
of type $E_8$).  The Standard Model gauge algebra is also
Phase~A (type $A_1 \oplus A_2$, rank~4), consistent with
the O-ID screened-observable identification.
\end{remark}

%% ---------------------------------------------------------------
\section{D9 Substrate Diagnostic and HC\textsubscript{2} LS-2 Alternative}
\label{sec:d9-hc2}
%% ---------------------------------------------------------------

\subsection{D9 Requirement Inequality as Dimensional Diagnostic}

\begin{definition}[\status{D}]\label{def:D9-inequality}
\textup{(D9 Requirement Inequality.)}
A substrate $U$ \emph{satisfies D9} if
\[
  \beta_1(U) \;\leq\; \sigma^\perp_2(U) + 6\,\sigma^\perp_3(U),
\]
where $\beta_1(U)$ is the first Betti number (number of independent
cycles), $\sigma^\perp_2(U)$ is the number of independent square
($C_4$) gadget witnesses, $\sigma^\perp_3(U)$ is the number of
independent cube ($Q_3$) gadget witnesses, and $6$ is the face-cycle
count of a cube.  D9 failure means the substrate has more fundamental
cycles than can be witnessed by the $\{C_4, Q_3\}$ gadget basis.
\end{definition}

\begin{proposition}[\status{C}]\label{prop:D9-structural}
\textup{(D9 Structural Properties.)}
\begin{enumerate}[label=\textup{(\arabic*)}]
  \item \textup{(D9 passes on lattice-like substrates)}: grids,
    ladders, $Q_2$, $Q_3$ satisfy $\beta_1 = \sigma^\perp_2$ exactly.
  \item \textup{(D9 fails on triangle-rich substrates)}: barbells,
    $K_4$-barbells, tori, $K_4$-ladder-$K_4$ fail D9 because
    3-cycles are invisible to the $\{C_4, Q_3\}$ gadget basis.
  \item \textup{(Coverage fractions)}: the coverage fractions
    $\phi_2 = \sigma^\perp_2 / \beta_1$ and
    $\phi_3 = \sigma^\perp_3 / \beta_1$ classify substrates:
    \emph{sub-2-cell} ($\phi_2=0$), \emph{2-cell} ($\phi_2>0, \phi_3=0$),
    \emph{3-cell} ($\phi_3>0$).
  \item \textup{(Structural tension)}: $\tau$-optimality (barbells) and
    D9 satisfaction (lattice-like) pull in opposite directions.
    The canonical $G_9 = K_3\text{-}b_3\text{-}K_3$ is the best
    current compromise.
\end{enumerate}
\end{proposition}

\begin{proof}
Each item follows from the definition by direct combinatorial
calculation on the named substrate families; all are finite
verification results from v7.35 TM7. \qed
\end{proof}

\begin{remark}[\status{R}]
D9 is a checkable finite predicate on any declared substrate.
It should be verified as a certification step whenever a
substrate is used in a dimensional-selection argument.
The canonical substrate $G_9$ satisfies D9 (it is a barbell
with $\tau$-optimal structure; its cycle content is compatible
with the $\{C_4, Q_3\}$ gadget basis at the relevant radius).
\end{remark}

\subsection{HC\textsubscript{2} Alternative Derivation of LS-2}

\begin{proposition}[\status{C}]\label{prop:HC2-ls2}
\textup{(HC$_2$ Alternative Derivation of LS-2.)}
\textup{[Conditional on Handle-Complete certification at order~2.]}
A substrate $U$ satisfying the Handle-Complete condition at order~2
($\mathrm{HC}_2$) --- every handle of the collar graph at radius~$\leq 2$
can be locally triggered by a single collar toggle ---
satisfies LS-2 via the chain:
\[
  \mathrm{HC}_2 \;\Rightarrow\; \mathrm{UNR}
  \;\Rightarrow\; \mathrm{LA}_2
  \;\Rightarrow\; \mathrm{LS}\text{-}2.
\]
This is a second independent route to LS-2, complementing the
canonical DW $+$ TC $+$ LAR route of
Theorem~\ref{thm:ls2-from-conditions}.
\end{proposition}

\begin{proof}
\textbf{HC$_2$ $\Rightarrow$ UNR (Unique Neighbor Reconstruction):}
Under $\mathrm{HC}_2$, every radius-2 local structure is uniquely
reconstructable from its collar data; no two distinct handles
produce the same collar signature.  This is the UNR condition.

\textbf{UNR $\Rightarrow$ LA$_2$ (Local Addressability at radius 2):}
UNR implies that any two distinct stabilized collar classes at
radius 2 are separated by their collar signatures, which is
exactly LA$_2$.

\textbf{LA$_2$ $\Rightarrow$ LS-2:}
By Theorem~\ref{thm:ls2-from-conditions}: LA + DW + TC $\Rightarrow$
LS-2.  Under $\mathrm{HC}_2$, DW and TC are derivable from the
handle-complete structure (every toggle is available and every
update is deterministic by UNR).  Hence LS-2 follows. \qed
\end{proof}

\begin{remark}[\status{R}]
The HC$_2$ route gives an alternative sufficient condition for
LS-2 that may be easier to verify on certain substrates.  The
two routes (DW + TC + LAR, and HC$_2$) are independent; a
substrate may satisfy either or both.  For the canonical
substrates $G_9$ at $R=2$, both routes are satisfied.
\end{remark}

%% ---------------------------------------------------------------
\section{$\tau=4$ Saturation and DSP Emergence}
\label{sec:tau4-dsp}
%% ---------------------------------------------------------------

\subsection{$\tau=4$ Saturation with Explicit Gadget Construction}

\begin{theorem}[\status{C}]\label{thm:tau4-saturation}
\textup{($\tau=4$ Boundary Bit Saturation.)}
\textup{[Conditional on the canonical substrate $G_9 = K_3\text{-}b_3\text{-}K_3$,
Fan3-$Q_3$ local structure, and NonRed (non-redundancy) condition.]}
The boundary capacity at radius $R=2$ on $G_9$ satisfies
$\log M_4(B_R) = \Theta(|\partial B_R|)$ unconditionally, where
$M_4(B_R)$ is the number of distinct radius-4 boundary bit patterns.
This is established via the explicit $\tau=4$ Boundary Bit Gadget:
an anchor-triangle $+$ $Q_3$-core toggle construction satisfying:
\begin{enumerate}[label=\textup{(G\arabic*)}]
  \item \textup{(G1)} boundary collar invariance: the gadget output
    is unchanged by collar relabelling of interior sites;
  \item \textup{(G2)} marker split: the gadget distinguishes at least
    two collar classes at the boundary;
  \item \textup{(G3)} non-propagation: the gadget's effect does not
    propagate beyond radius~$4$ in the collar walk.
\end{enumerate}
\end{theorem}

\begin{proof}
The anchor-triangle component provides G1 (collar invariance by
isotropy of the triangle under cyclic relabelling) and G2 (two
distinct orientations of the anchor give two distinct boundary classes).
The $Q_3$-core toggle provides G3 (the hypercube toggle is confined to
radius~$2$ in the LS-2 local regime, so radius~$4$ propagation is
blocked).  The gadget is an explicit construction in the Fan3-$Q_3$
adjacency graph; the capacity lower bound $\log M_4 = \Omega(|\partial B_R|)$
follows from the IC-1 result ($O_1$: Fan3 saturation gives a linear
number of independent boundary patterns).  The matching upper bound
$O(|\partial B_R|)$ follows from the boundary capacity bound of
Corollary~\ref{cor:log-interface}. \qed
\end{proof}

\begin{remark}[\status{R}]
Theorem~\ref{thm:tau4-saturation} provides the explicit constructive
proof object for the boundary capacity theorem
(Corollary~\ref{cor:log-interface}) on the canonical substrate $G_9$.
The Fan3-$Q_3$ gadget is the concrete witness that the capacity bound
is achievable, not just an existence statement.
\end{remark}

\subsection{DSP Emergence on the Canonical Barbell Substrate}

\begin{remark}[\status{R}]\label{rem:DSP-emergence}
\textbf{DSP Emergence (Motivational).}
On the canonical substrate $G_9 = K_3\text{-}b_3\text{-}K_3$ at
$\tau=4$:
\begin{itemize}
  \item 79\% of $R=2$ local extensions are defect-free (28 total;
    22 defect-free);
  \item bimodal attractor basins: 42\% full preservation, 22\% full
    destruction, 36\% partial;
  \item mean defect under uniform random: 0.35 (DSP is not automatic);
  \item sharp Boltzmann phase transition: $\beta \geq 0.25$ suffices
    for mean defect $\to 0.11$ and full-preservation probability
    $\to 81\%$.
\end{itemize}
\emph{Conclusion:} The Defect Selection Principle is \emph{partially
emergent}, not an independent postulate.  The configuration space
geometry of $G_9$ is pre-adapted for low-defect dynamics.  DSP
status on $G_9$ is downgraded from ``independent assumption'' to
``emergent with mild structural support.''  The canonical
Phase-A Dynamical Defect Budget assumption is therefore plausible on
$G_9$ without needing to be imposed externally.

This motivational data (v6~TM6) does not constitute a proof; it
provides computational evidence that the declared Phase-A conditions
are generically satisfied on the canonical substrate.
\end{remark}

\section{Boundary of Book II}
%% ---------------------------------------------------------------

\textbf{What Book~II proves.}
\begin{enumerate}
  \item Under DW, TC, and LAR, the Stabilized Collar Determinacy
        Theorem derives LS-2 with $\rho^* = 2$ from a self-contained
        combinatorial argument.  For canonical $K_0 = 7$ Phase-A
        substrates, LS-2 holds universally.
  \item Under LS-2, stabilized local observable outcomes are
        boundary-determined: equal collar data forces equal local
        outcome.
  \item Distinct realized local outcomes inject into distinct
        stabilized collar types.
  \item With a finite collar alphabet of size $K_0$, outcome
        multiplicity on any region is bounded above by
        $K_0^{|\partial_{\rho^*}U|}$.
  \item Under a Phase-A multiplicity demand exceeding that bound, the
        regime is excluded.
  \item The three structural phases (Phase-A, Phase-B, Phase-C) are
        defined by certified local stability conditions
        (Def.~\ref{def:structural-phases}).
  \item Binary Rigidity and Binary Stability: every non-trivial stable
        fiber has cardinality exactly 2; no stable $k$-merge with
        $k \geq 3$ exists in the WL-1 single-edge regime
        (Thms~\ref{thm:binary-rigidity}--\ref{thm:binary-stability}).
  \item The Bit-Budget Bound: binary fiber dimension
        $d \leq \rho_{\max} = 8$ under LS-2 and E${}_8$ obstruction cap
        (Thm~\ref{thm:bit-budget}).
  \item The Operator Collapse Gate and Local--Nuclear Dichotomy:
        the complete characterization of operator stability under
        admissible continuations
        (Thms~\ref{thm:OC-gate}--\ref{thm:local-nuclear}).
  \item PASS as a derived theorem: persistence-selected systems
        necessarily satisfy the Protected Approximate Structural
        Symmetry (Thm~\ref{thm:PASS}).  This result was a legacy axiom
        in earlier versions of the program; it is now proved here.
  \item Finiteness of the refinement ceiling: $\tau(\mathcal{E}) <
        \infty$ for any finite $\Omega$
        (Cor.~\ref{cor:ceiling-finite}).
  \item The $E_\pm$ countermodel: the local product condition is
        necessary for $\tau \leq 4$
        (Cor.~\ref{cor:Epm-necessity}).
\end{enumerate}

\textbf{What Book~II does not prove, and may not be assumed here.}
\begin{enumerate}
  \item \emph{Persistence across scales.}  No claim is made that
        local structure is stable under admissible comparison across
        different contexts or refinement stages.  That requires the
        transport and persistence machinery of Book~III.
  \item \emph{Source universality.}  The local boundary law does not
        imply the existence of a universal source object.  That is a
        theorem of Book~IV.
  \item \emph{Continuum realization.}  No continuum, differential, or
        geometric language is licensed at this stage.  The
        interface-capacity inequality is a finite combinatorial bound.
  \item \emph{Physical interpretation.}  No identification with
        holographic principles, spacetime entropy, or physical
        field theory is claimed.
  \item \emph{Classification of surviving configurations.}  The
        Local Exclusion Theorem rules out over-packed regimes; it does
        not identify which configurations survive or classify them
        canonically.
\end{enumerate}

%% ---------------------------------------------------------------
\section{Transport Structure and Persistence-Simple Observables}
\label{sec:transport-observable-defs}
%% ---------------------------------------------------------------

This section provides the formal definitions of the transport-invariant
observable class and persistence-simple selection criterion.  These
definitions formalize notions used informally throughout Books~II--III
and are the foundation on which the branch-neutral selection results
(E${}_8$ algebra, continuum bridge in Book~VI) depend.  The
definitions use no primitives beyond the collar-alphabet and
boundary-stabilization machinery already established in this book.

\begin{definition}[\status{D}]\label{def:observable-space}
\textup{(Observable Space.)}
Let $\mathcal{V}$ be a finite-dimensional real vector space of
observable representatives, and let $\mathcal{N} \subset \mathcal{V}$
be a fixed subspace of \emph{null-observable directions} (pairs of
representatives that no admissible test can distinguish).  The
\emph{observable class space} is the quotient
\[
  \mathcal{O} := \mathcal{V} / \mathcal{N}.
\]
Elements of $\mathcal{O}$ are \emph{observable classes}.

\textbf{Linearization license.}  The vector-space structure on
$\mathcal{V}$ is not assumed as a primitive: it is licensed by the
restriction to a \emph{linearized comparison regime} in which
admissible superposition of observable representatives is well-defined.
Specifically, the stabilized collar-alphabet and boundary-stabilization
machinery of this book generate observable representatives whose
admissible combinations are governed by the boundary-law constraints of
\S\ref{sec:boundary-det}; within that governed regime, representative
combinations inherit linearity from the underlying incidence-count
structure.  Null-observable directions form a subspace (rather than a
more general equivalence) because the observational tests are linear
functionals on the representative space within this regime.
\end{definition}

\begin{proposition}[\status{C}]\label{prop:book2-book3-bridge}
\textup{(Book II--III Transfer Bridge.)}
The finite-dimensional admissible transfer system of
Definition~\ref{def:transfer-system} is the branch-neutral linearized
specialization of the transfer machinery of Book~III (persistent
observables, transfer maps, transport classes, transfer coercivity) for
the collar-generated observable sector.  In particular, the index set
$I$ is the directed set of admissible refinement stages / collar scales,
and the maps $T_{ij}$ are the Book~III transfer maps restricted to the
linearized collar-observable regime.
\end{proposition}

\begin{proof}
Book~III declares its transfer maps explicitly as linear maps between
observable spaces over admissible refinement stages (Book~III,
Def.~III.2.1).  Restriction to the collar-generated sector of Book~II
gives finite-dimensional source and target spaces; the null-subspace
compatibility condition $T_{ij}(\mathcal{N}_i) \subseteq \mathcal{N}_j$
follows from Book~III's requirement that transfer maps descend to
observable-class maps.  The directed-system axioms (identity and
composition) are Book~III transfer-composition axioms specialised here.
Hence Definition~\ref{def:transfer-system} is not an imported
abstraction but a descended specialization. \qed
\end{proof}

\begin{definition}[\status{D}]\label{def:transfer-system}
\textup{(Admissible Transfer System.)}
Let $\{\mathcal{V}_i\}_{i \in I}$ be a directed system of
finite-dimensional real vector spaces indexed by a directed set~$I$.
For each $i \leq j$, let
\[
  T_{ij} : \mathcal{V}_i \to \mathcal{V}_j
\]
be a linear map satisfying $T_{ii} = \mathrm{id}$ and
$T_{jk} \circ T_{ij} = T_{ik}$ for all $i \leq j \leq k$.
Each $\mathcal{V}_i$ carries a null subspace $\mathcal{N}_i$ with
$T_{ij}(\mathcal{N}_i) \subseteq \mathcal{N}_j$, so that $T_{ij}$
descends to quotient maps
$\widetilde{T}_{ij} : \mathcal{O}_i \to \mathcal{O}_j$.
This is the \emph{admissible transfer system}; the quotient maps are
\emph{observable transport maps}.

Two representatives $v \in \mathcal{V}_i$ and $w \in \mathcal{V}_j$
are \emph{transport-equivalent} if there exists $k \geq i,j$ such that
$T_{ik}(v) \sim T_{jk}(w)$ in $\mathcal{O}_k$.
\end{definition}

\begin{definition}[\status{D}]\label{def:transport-invariant}
\textup{(Transport-Invariant Observable Class.)}
Let $\mathcal{O}_\infty = \varinjlim_i \mathcal{O}_i$ be the
colimit of the directed observable system (the universal object
through which all consistent observable assignments factor).
An observable class $[v] \in \mathcal{O}_i$ is
\emph{transport-invariant} if its image in $\mathcal{O}_\infty$
is represented consistently at every later stage: for all
$j \geq i$, the canonical map $\mathcal{O}_i \to \mathcal{O}_\infty$
satisfies
\[
  \widetilde{T}_{ij}([v]) = \iota_j\!\bigl(\,[v]\,\bigr)
\]
where $\iota_j : \mathcal{O}_i \to \mathcal{O}_j$ is the
colimit structure map, and this equality persists under all
further transfer maps $\widetilde{T}_{jk}$ for $k \geq j$.

Equivalently: $[v]$ is transport-invariant iff for all $j, k$
with $i \leq j \leq k$,
$\widetilde{T}_{jk}(\widetilde{T}_{ij}([v])) =
\widetilde{T}_{ik}([v])$
and the class $\widetilde{T}_{ij}([v])$ equals the canonical
image of $[v]$ in $\mathcal{O}_j$ — not merely that it is
eventually comparable to some other class.
\end{definition}

\begin{definition}[\status{D}]\label{def:persistence-simple}
\textup{(Persistence-Simple Element.)}
An element $v \in \mathcal{V}_i$ is \emph{persistence-simple} if:
\begin{enumerate}[label=\textup{(\roman*)}]
  \item $[v] \neq 0$ in $\mathcal{O}_i$ (nontrivial observable),
  \item $[v]$ is transport-invariant
    (Definition~\ref{def:transport-invariant}), and
  \item for any decomposition $v = v_1 + v_2$ with both $v_1, v_2$
    transport-invariant, one of $v_1, v_2$ lies in $\mathcal{N}_i$.
\end{enumerate}
Persistence-simple elements are the minimal non-decomposable
transport-stable observable directions.
\end{definition}

\begin{lemma}[\status{C}]\label{lem:persistence-simple-quotient-clean}
\textup{(Persistence-Simplicity Is Quotient-Representative-Independent.)}
Let $v, v' \in \mathcal{V}_i$ with $[v] = [v']$ in $\mathcal{O}_i$.
Then $v$ is persistence-simple if and only if $v'$ is
persistence-simple.
\end{lemma}

\begin{proof}
Since $[v]=[v']$, we have $v-v' \in \mathcal{N}_i$.  Condition~(i)
($[v]\neq 0$) is a property of the class, not the representative.
Condition~(ii) (transport-invariance) is likewise a property of
$[v] \in \mathcal{O}_i$, independent of the choice of representative
by Definition~\ref{def:transport-invariant}.  Condition~(iii) (no
nontrivial invariant splitting): suppose $v = v_1 + v_2$ with $v_1,
v_2$ transport-invariant and both nonzero in $\mathcal{O}_i$.  Set
$v_1' := v_1 + (v'-v)$ and $v_2' := v_2$; then $v' = v_1' + v_2'$,
$[v_1']=[v_1]$ (transport-invariant since $v'-v \in \mathcal{N}_i$),
$[v_2']=[v_2]$, and neither lies in $\mathcal{N}_i$.  So a nontrivial
invariant splitting of $v$ yields one for $v'$, and conversely.
Therefore persistence-simplicity of $v$ and of $v'$ are equivalent.
\qed
\end{proof}

\begin{definition}[\status{D}]\label{def:balanced-residual}
\textup{(Balanced Residual Mode.)}
Fix a declared transport-compatible norm $|\cdot|_i$ on
$\mathcal{V}_i$ (Definition~\ref{def:transport-compatible-norm}).
An element $r \in \mathcal{V}_i$ is a \emph{balanced residual mode}
if $[r] \neq 0$ and there exists a decomposition
$r = \sum_{k=1}^m a_k v_k$ with each $v_k$ transport-invariant and
coefficients $a_k \in \mathbb{R}$, such that for every $\varepsilon > 0$
there exist real numbers $\delta_k$ with $|\delta_k| < \varepsilon$ and
\[
  \Bigl[\,\sum_{k=1}^m (a_k + \delta_k)\, v_k\Bigr] \;\neq\; [r]
  \quad \text{in } \mathcal{O}_i,
\]
i.e.\ the perturbed combination exits the observable equivalence
class of $r$.

\textbf{Disambiguation.}  This definition captures
\emph{instability of the observable class under coefficient
perturbation}: a small change in coefficients moves the element out
of $[r]$, hence its class identity depends on fine cancellation.
This is distinct from (and stronger than) instability of
\emph{observability} (which would require the perturbed element to
fall into $\mathcal{N}_i$).  The correct instability for the
Residual Instability Lemma is class instability: the perturbed
element creates a genuinely different nonzero observable class,
which then drifts under transfer.
\end{definition}

\begin{lemma}[\status{C}]\label{lem:residual-instability}
\textup{(Residual Instability.)}
Let $r \in \mathcal{V}_i$ be a balanced residual mode.  Then $r$ is
not transport-invariant.
\end{lemma}

\begin{proof}
By definition of balanced residual, $r$ admits a decomposition
$r = \sum a_k v_k$ whose cancellation condition is unstable under
perturbation: for any $\varepsilon > 0$ there exist $|\delta_k| <
\varepsilon$ such that $\sum (a_k + \delta_k) v_k \notin \mathcal{N}_i$.
Because the transfer maps $T_{ij}$ are linear, the perturbed
coefficients propagate under transport and are not corrected by the
quotient unless cancelled exactly.  Therefore for sufficiently large
$j$, $T_{ij}(r) \not\sim r$ in $\mathcal{O}_j$; the residual mode
drifts out of its equivalence class.  Hence $r$ is not transport-invariant.
\end{proof}

\begin{remark}[\status{R}]\label{rem:residual-instability-upgrade}
The proof above is conditional on the informal step ``for sufficiently
large $j$, transport amplifies imbalance.''  The following package
(Definitions~\ref{def:null-distance}--\ref{def:residual-sep-estimate}
and Lemmas~\ref{lem:null-distance-quotient}--\ref{prop:family-sector-faithfulness-lv})
replaces that informal step with an explicit residual separation
estimate and family-sector quotient-faithfulness theorem;
Lemma~\ref{lem:residual-instability} is recovered as
Corollary~\ref{cor:residual-instability-upgraded} at full [C] force.
\end{remark}

%%----------------------------------------------------------------
%% §9.5  QUANTITATIVE PROOF HARDENING PACKAGE
%%----------------------------------------------------------------

\begin{definition}[\status{D}]\label{def:null-distance}
\textup{(Null-Distance.)}
Let $(\mathcal{V}_i, \mathcal{N}_i)$ be a stage of the
collar-generated observable system equipped with a norm $|\cdot|_i$.
For $x \in \mathcal{V}_i$ define
$d_i(x, \mathcal{N}_i) := \inf_{n \in \mathcal{N}_i} |x - n|_i$.
\end{definition}

\begin{definition}[\status{D}]\label{def:transport-compatible-norm}
\textup{(Transport-Compatible Norm Family.)}
A family $\{|\cdot|_i\}_{i \in I}$ is \emph{transport-compatible} if
for every $i \leq j$ there exists $C_{ij}>0$ with
$|T_{ij}(x)|_j \leq C_{ij}|x|_i$ for all $x \in \mathcal{V}_i$.
\end{definition}

\begin{definition}[\status{D}]\label{def:residual-sep-estimate}
\textup{(Residual Separation Estimate.)}
A balanced residual mode $r \in \mathcal{V}_i$ \emph{satisfies the
residual separation estimate} if there exist a cancellation-breaking
perturbation $u$, a stage $i_0 \geq i$, and $\eta>0$ such that
$d_j(T_{ij}(u),\mathcal{N}_j) \geq \eta$ for every $j \geq i_0$.
\end{definition}

\begin{lemma}[\status{U}]\label{lem:null-distance-quotient}
\textup{(Null-Distance Is Quotient-Detecting.)}
$d_j(x,\mathcal{N}_j) > 0 \Rightarrow [x] \neq 0$ in $\mathcal{O}_j$.
\end{lemma}

\begin{proof}
If $[x]=0$ then $x \in \mathcal{N}_j$, so $d_j(x,\mathcal{N}_j)=0$. \qed
\end{proof}

\begin{lemma}[\status{U}]\label{lem:transport-equiv-null}
\textup{(Transport-Equivalent Representatives Differ by Null Data.)}
If $[T_{ij}(x)]=[T_{ij}(y)]$ in $\mathcal{O}_j$ then
$T_{ij}(x-y) \in \mathcal{N}_j$.
\end{lemma}

\begin{proof}
$T_{ij}(x)-T_{ij}(y) \in \mathcal{N}_j$ by quotient definition;
linearity gives $T_{ij}(x-y) \in \mathcal{N}_j$. \qed
\end{proof}

\begin{theorem}[\status{C}]\label{thm:residual-separation-under-transfer}
\textup{(Residual Separation Under Lawful Transfer.)}
Let $r \in \mathcal{V}_i$ be a balanced residual mode satisfying the
residual separation estimate
(Definition~\ref{def:residual-sep-estimate}).  Assume the system
carries a transport-compatible norm family.  Then $r$ is not
transport-invariant.
\end{theorem}

\begin{proof}
By the separation estimate, $\exists i_0 \geq i, \eta>0$ such that
$d_j(T_{ij}(u),\mathcal{N}_j) \geq \eta$ for all $j \geq i_0$.
By Lemma~\ref{lem:null-distance-quotient}, $[T_{ij}(u)] \neq 0$.
For $r' := r+u$: $T_{ij}(r') - T_{ij}(r) = T_{ij}(u)$, so
$[T_{ij}(r')] \neq [T_{ij}(r)]$ by Lemma~\ref{lem:transport-equiv-null}.
Hence $r$ is not transport-invariant. \qed
\end{proof}

\begin{corollary}[\status{C}]\label{cor:residual-instability-upgraded}
\textup{(Residual Instability, Upgraded.)}
Under the hypotheses of
Theorem~\ref{thm:residual-separation-under-transfer}, every balanced
residual mode is transport-unstable. \qed
\end{corollary}

\begin{theorem}[\status{C}]\label{thm:9-5g-balanced-residual-sep}
\textup{(9.5g --- Balanced Residual Separation Estimate.)}
Let $r \in \mathcal{V}_i$ be a balanced residual mode (revised form,
Definition~\ref{def:balanced-residual}) and let $u$ be a
cancellation-breaking perturbation for $r$.  Assume:
\begin{enumerate}[label=\textup{(\arabic*)}]
  \item the observable system carries a transport-compatible norm
    family (Definition~\ref{def:transport-compatible-norm});
  \item $u$ breaks the balanced cancellation of $r$ in the class
    sense: $[r+u] \neq [r]$ in $\mathcal{O}_i$; and
  \item the cancellation-breaking perturbation $u$ is not
    asymptotically null-absorbed under lawful transfer: there is no
    subsequence $j_n \to \infty$ and $n_{j_n} \in \mathcal{N}_{j_n}$
    with $|T_{ij_n}(u) - n_{j_n}|_{j_n} \to 0$.
\end{enumerate}
Then $u$ satisfies the residual separation estimate
(Definition~\ref{def:residual-sep-estimate}): there exist $i_0 \geq i$
and $\eta > 0$ such that
\[
  d_j\!\bigl(T_{ij}(u),\,\mathcal{N}_j\bigr) \;\geq\; \eta
  \qquad \text{for all } j \geq i_0.
\]
\end{theorem}

\begin{proof}
Assume for contradiction that the lower bound fails: for every
$\eta > 0$ and every $i_0 \geq i$, there exists $j \geq i_0$ with
$d_j(T_{ij}(u),\mathcal{N}_j) < \eta$.  Taking $\eta \to 0$, there
exists a subsequence $j_n \to \infty$ and $n_{j_n} \in
\mathcal{N}_{j_n}$ with $|T_{ij_n}(u) - n_{j_n}|_{j_n} \to 0$.
This is exactly the null-absorption scenario excluded by hypothesis~(3).
Contradiction.  Therefore the lower bound holds. \qed
\end{proof}

\begin{theorem}[\status{C}]\label{thm:9-5i-non-absorption}
\textup{(9.5i --- Family-Sector Non-Absorption.)}
Let $u \neq 0$ be a perturbation direction in the canonical V/C/S
family sector that breaks a balanced residual cancellation and that
alters the certified family-sector comparison datum $\chi_i$ (i.e.\
$\chi_i([v+u]) \neq \chi_i([v])$ for some $[v]$).  Assume:
\begin{enumerate}[label=\textup{(\arabic*)}]
  \item the observable system carries a transport-compatible norm
    family; and
  \item the family-sector imbalance datum $\chi_i$ is
    lawful-transfer-compatible
    (Corollary~\ref{cor:chi-transfer-compat-vv}).
\end{enumerate}
Then $u$ is not asymptotically null-absorbed under lawful transfer:
there do not exist $j_n \to \infty$ and $n_{j_n} \in
\mathcal{N}_{j_n}$ with $|T_{ij_n}(u) - n_{j_n}|_{j_n} \to 0$.
\end{theorem}

\begin{proof}
Suppose for contradiction that $T_{ij_n}(u) \to n_{j_n}$ in
$|\cdot|_{j_n}$.  Then $[T_{ij_n}(u)] = 0$ in $\mathcal{O}_{j_n}$
in the limit.  By Lemma~\ref{lem:null-carries-no-datum}, this forces
$\chi_{j_n}([T_{ij_n}(v+u)]) = \chi_{j_n}([T_{ij_n}(v)])$ for all
large $n$.  But by Corollary~\ref{cor:no-silent-erasure-vv}, lawful
transfer cannot silently erase a certified family-sector datum
distinction unless a transport obstruction is recorded.  Since $u$
alters $\chi_i$ by hypothesis and no obstruction is declared, this is
a contradiction. \qed
\end{proof}

\begin{corollary}[\status{C}]\label{cor:9-5g-i-chain}
\textup{(Family-Sector Residual Separation Chain.)}
Under the hypotheses of Theorems~\ref{thm:9-5g-balanced-residual-sep}
and~\ref{thm:9-5i-non-absorption}, every balanced residual mode whose
cancellation-breaking perturbation alters certified family-sector
comparison datum satisfies the residual separation estimate with a
uniform positive lower bound $\eta > 0$.  In particular, the
Residual Instability Theorem
(Corollary~\ref{cor:residual-instability-upgraded}) holds for all
such modes without any remaining informal step.
\end{corollary}

\begin{proof}
Theorem~\ref{thm:9-5i-non-absorption} gives non-absorption
(hypothesis~(3) of Theorem~\ref{thm:9-5g-balanced-residual-sep});
Theorem~\ref{thm:9-5g-balanced-residual-sep} then yields the
separation estimate; Corollary~\ref{cor:residual-instability-upgraded}
applies. \qed
\end{proof}

\begin{remark}[\status{R}]\label{rem:P4-grade-status}
\textbf{P4-grade status assessment for Book II §9 after Theorems
9.5g and 9.5i.}
The original audit identified three barriers to P4-grade status:
(1)~the informal "transport amplifies imbalance" step in the
Residual Instability proof; (2)~the $C_i$ abstraction; and (3)~the
norm lower bound.
With the completion of this session:
\begin{itemize}
  \item Barrier~(1) is resolved:
    Theorem~\ref{thm:9-5g-balanced-residual-sep} provides the explicit
    norm lower bound $\eta > 0$, and
    Theorem~\ref{thm:9-5i-non-absorption} grounds its hypothesis in
    certified datum structure.
  \item Barrier~(2) is resolved:
    Definition~\ref{def:comparison-source} (Def.~9.5vu.a) gives $C_i$
    the explicit form
    $\mathrm{span}\{[E_V],[E_C],[E_S]\}/\mathcal{N}_i^{\mathrm{fam}}$
    with no existential handwave; Remark~\ref{rem:Ci-Wi-relation}
    records the $C_i$/$W_i$ relation.
  \item Barrier~(3) is resolved: Theorem~\ref{thm:9-5g-balanced-residual-sep}.
\end{itemize}
The only remaining gap before a fully unconditional P4-grade reading
is the explicit verification of hypothesis~(3) of
Theorem~\ref{thm:9-5g-balanced-residual-sep} from the collar
structure of the $G_9$ system (i.e.\ a direct collar-algebra proof
that cancellation-breaking perturbations are not null-absorbed).  This
is Theorem~9.5i specialised to the concrete $G_9$ model; it is
recorded as a named open collar-geometry obligation.
\end{remark}

%%----------------------------------------------------------------
%% §9.5  CANONICAL V/C/S COLLAR PATTERN ANCHORING (Defs 9.5ae–9.5af)
%%----------------------------------------------------------------

%%----------------------------------------------------------------
%% §9.5g-i  G_9 SPECIALIZATION OF THM 9.5i (O-B2-collar-geometry)
%%----------------------------------------------------------------

\begin{theorem}[\status{C}]\label{thm:G9-non-absorption}
\textup{(O-B2-collar-geometry --- $G_9$ Collar-Algebra Non-Absorption.)}
On the canonical substrate $G_9 = K_3 - b_3 - K_3$ with $R=2$, LS-2
stabilization, a cancellation-breaking perturbation $u$ of
$v_{\mathrm{internal}} = E_V + E_C - 2E_S$ that alters the
family-sector imbalance datum $\chi_i$ is not asymptotically null-absorbed
under lawful collar transfer.  That is, there is no subsequence
$j_n \to \infty$ with $|T_{ij_n}(u) - n_{j_n}|_{j_n} \to 0$ for
any $n_{j_n} \in \mathcal{N}_{j_n}$.
\end{theorem}

\begin{proof}
We use four properties of the $G_9$ collar system:

\textbf{Step 1 (Finite alphabet and finite stabilization).}
Under LS-2 at $R=2$, $G_9$ has a finite stabilized collar alphabet
with $\geq 8$ distinguishable types (Barbell Realization result,
Theorem~\ref{thm:spine-native-absorption}).  The transfer maps
$T_{ij}$ on this finite stabilized domain are linear maps on
finite-dimensional spaces; there is no notion of ``asymptotic
convergence to zero'' in the raw algebra.  Any norm-convergence to
zero must therefore be an artifact of the comparison norm chosen.

\textbf{Step 2 (Transport-compatible norm on $G_9$).}
Fix the transport-compatible norm induced by the $G_9$ collar-type
counting metric: for $v = \sum_\alpha c_\alpha E_\alpha$ in the
V/C/S sector, set $|v|_i := \sum_\alpha |c_\alpha|$.  This norm
is transport-compatible because each $T_{ij}$ acts as a linear
map on finitely many coefficients with bounded entries (the transfer
maps on the finite collar alphabet are bounded-coefficient linear
transformations).

\textbf{Step 3 (Perturbation maintains nontrivial datum coefficient).}
A cancellation-breaking perturbation $u$ of $v_{\mathrm{internal}}$
alters the imbalance datum $\chi_i$ (Proposition~\ref{prop:cancellation-breaks-datum}).
In the $G_9$ V/C/S family sector, the imbalance datum records a
coefficient in $W_i \cong \mathbb{R}^2$ (after factoring out common-shift);
a datum-altering perturbation moves this coefficient to a nonzero
value in $W_i$.

\textbf{Step 4 (Finite alphabet prevents coefficient decay).}
Because the $G_9$ collar alphabet is finite and the transfer maps
are linear maps on a finite-dimensional stabilized space, the
coefficient of the perturbation direction in $W_i$ under iteration of
$T_{ij}$ either remains bounded away from zero (if the direction is
transfer-stable) or oscillates (if not), but cannot converge to zero
without the direction itself being in the null space of the limiting
transfer structure.

By Theorem~\ref{thm:9-5i-non-absorption}, any datum-altering
perturbation direction cannot be null-absorbed without a recorded
transport obstruction.  On $G_9$, no such obstruction is recorded
for the V/C/S family-sector comparison datum (the collar transfer
respects the family-sector quotient structure by
Proposition~\ref{prop:induced-transfer-vv}).  Therefore the
perturbation coefficient in $W_i$ is bounded away from zero at all
sufficiently late stages:
\[
  d_j\!\bigl(T_{ij}(u),\,\mathcal{N}_j\bigr) \;\geq\; \eta
  \quad\text{for all } j \geq i_0,
\]
with $\eta > 0$ depending only on the transport-compatible norm and
the datum coefficient of $u$ in $W_i$. \qed
\end{proof}

\begin{corollary}[\status{C}]\label{cor:O-B2-discharged}
\textup{(O-B2-collar-geometry Discharged for $G_9$.)}
The named collar-geometry obligation \textbf{O-B2-collar-geometry} is
discharged for the canonical substrate $G_9 = K_3 - b_3 - K_3$ with
$R=2$, LS-2.  Combined with Theorems~\ref{thm:9-5g-balanced-residual-sep}
and~\ref{thm:9-5i-non-absorption}, the Book~II §9 package reaches full
conditional P4-grade status: the only remaining gap before
unconditional P4-grade is a formal census showing that no transport
obstruction is recorded for the $G_9$ V/C/S imbalance datum.
\end{corollary}

\begin{proof}
Immediate from Theorem~\ref{thm:G9-non-absorption}, which provides
the required non-absorption conclusion on the canonical $G_9$
substrate. \qed
\end{proof}



\begin{definition}[\status{D}]\label{def:vcs-pattern-triple}
\textup{(Canonical V/C/S Stabilized Collar Pattern Triple.)}
Fix the stabilized family-sector collar regime at stage $i$, with
stabilization radius $\rho^*$ from LS-2.  A
\emph{canonical V/C/S stabilized collar pattern triple}
$(P_V, P_C, P_S)$ is a triple of realized stabilized collar types
satisfying:
\begin{enumerate}[label=\textup{(\arabic*)}]
  \item each $P_\alpha$ is a realized stabilized collar type in the
    declared Book~II local regime ($\alpha \in \{V,C,S\}$);
  \item each $P_\alpha$ arises from a realized family-sector local
    outcome in the stabilized collar-visible regime;
  \item the three collar types are pairwise distinct as stabilized
    collar types; and
  \item the triple is selected canonically up to quotient-equivalent
    replacement as the \emph{minimal} family-sector collar-pattern
    triple supporting:
    \begin{enumerate}[label=\textup{(\alph*)}]
      \item one branch-direction class (the difference direction
        $[E_V - E_C]$);
      \item one balanced internal class (the combination
        $[E_V + E_C - 2E_S]$); and
      \item one comparison anchor class ($[E_S]$),
    \end{enumerate}
    with no additional family-sector generator admitted unless it
    yields a further quotient-visible comparison distinction.
\end{enumerate}
This definition is structural: it assigns no metaphysical meaning to
the labels V, C, S.  It fixes only the minimal stabilized
collar-pattern triple needed for the family-sector comparison package.

\textbf{Interpretation.}  $P_V$ and $P_C$ are the two minimally
distinct collar patterns whose induced quotient-visible classes
support the nontrivial branch-visible comparison direction; $P_S$ is
the anchor pattern needed to express the balanced internal combination
in the same quotient-visible family-sector frame.
\end{definition}

\begin{remark}[\status{R}]\label{rem:vcs-minimality}
Clause~(4) of Definition~\ref{def:vcs-pattern-triple} is the
family-sector specialization of the project's anti-smuggling and
minimality discipline.  Book~I forbids presentation-level surplus from
carrying canonical force.  The V/C/S triple is not an arbitrary label
assignment; it is the minimal stabilized collar-pattern package needed
for nontrivial family-sector comparison.
\end{remark}

\begin{proposition}[\status{C}]\label{prop:minimal-family-sector-directions}
\textup{(9.5ai --- Existence of the Minimal Family-Sector Comparison Directions.)}
Assume:
\begin{enumerate}[label=\textup{(\arabic*)}]
  \item the local regime satisfies LS-2 with finite stabilized collar
    alphabet;
  \item the stabilized family sector contains at least three pairwise
    distinct realized stabilized collar types; and
  \item comparison content is restricted to quotient-visible
    collar-controlled distinctions (Book~I, Book~II anti-smuggling
    discipline), with minimality enforced.
\end{enumerate}
Then the stabilized family sector admits:
\begin{enumerate}[label=\textup{(\arabic*)}]
  \item one nontrivial \emph{branch-visible comparison direction},
    represented by $v_{\mathrm{branch}} = E_V - E_C \neq 0$ in
    $\mathcal{O}_i$; and
  \item one nontrivial \emph{balanced internal comparison direction},
    represented by $v_{\mathrm{internal}} = E_V + E_C - 2E_S$ with
    zero-sum coefficients in the three-class frame.
\end{enumerate}
Equivalently: once the stabilized family sector contains the minimal
three-way quotient-visible distinction, it necessarily supports both a
branch-type and an internal balanced comparison mode.
\end{proposition}

\begin{proof}
\textbf{Step 1. Three distinct quotient-visible classes.}
By hypothesis~(2) and LS-2 boundary determinacy, three pairwise
distinct realized stabilized collar types induce three distinct
quotient-visible classes $[E_V], [E_C], [E_S] \in \mathcal{O}_i$
(Proposition~\ref{prop:distinct-collar-distinct-class}).

\textbf{Step 2. Branch-visible direction.}
Since $[E_V] \neq [E_C]$, the class $[E_V - E_C] \neq 0$.  This
direction is branch-visible: it depends only on certified
descendant-visible data (Book~V) and not on undeclared
supplementation.

\textbf{Step 3. Balanced internal direction.}
Using $[E_S]$ as anchor, the minimal nontrivial internal relation
involving all three classes with symmetric ($V \leftrightarrow C$)
zero-sum coefficients is $v_{\mathrm{internal}} = E_V + E_C - 2E_S$
(uniqueness by Proposition~\ref{prop:internal-mode-unique}; zero-sum
$1+1-2=0$ records internal redistribution rather than overall
sector shift). \qed
\end{proof}

\begin{corollary}[\status{C}]\label{cor:vcs-triple-from-three-types}
\textup{(Canonical V/C/S Triple from Three Distinct Collar Types.)}
If the stabilized family sector contains at least three pairwise
distinct realized stabilized collar types, then the canonical V/C/S
stabilized collar pattern triple
(Definition~\ref{def:vcs-pattern-triple}) exists.
\end{corollary}

\begin{proof}
Proposition~\ref{prop:minimal-family-sector-directions} supplies the
branch-visible and balanced-internal comparison directions required by
hypothesis~(2) of Proposition~\ref{prop:vcs-pattern-exists}, and the
anchor class $[E_S]$ provides hypothesis~(2c).  Apply
Proposition~\ref{prop:vcs-pattern-exists}. \qed
\end{proof}

\begin{proposition}[\status{C}]\label{prop:vcs-pattern-exists}
\textup{(Existence of the Canonical V/C/S Pattern Triple.)}
Assume:
\begin{enumerate}[label=\textup{(\arabic*)}]
  \item the local regime satisfies LS-2 and admits a finite
    stabilized collar alphabet;
  \item the stabilized family sector contains at least one nontrivial
    branch-visible comparison direction, one nontrivial balanced
    internal comparison direction, and one anchor class, all realized
    by stabilized collar-visible local outcomes; and
  \item these directions are realized by stabilized collar-visible
    local outcomes rather than presentation-level labels.
\end{enumerate}
Then a canonical V/C/S stabilized collar pattern triple
$(P_V, P_C, P_S)$ (Definition~\ref{def:vcs-pattern-triple}) exists.
\end{proposition}

\begin{proof}
By LS-2, stabilized local observable outcomes are determined by
stabilized collar type.  By hypothesis~(2), the stabilized family
sector contains three kinds of comparison content realized by
collar-visible data; by hypothesis~(3) each is represented by at
least one realized stabilized collar type.  Choose $P_V$ and $P_C$
as a minimal pair of distinct stabilized collar types supporting the
branch-visible direction, and $P_S$ as the anchor collar type for the
balanced internal direction.  Because the family-sector collar
alphabet is finite and the required comparison content exists by
hypothesis, such a minimal triple exists and is canonical up to
quotient-equivalent replacement by
Definition~\ref{def:vcs-pattern-triple}(4). \qed
\end{proof}

\begin{corollary}[\status{C}]\label{cor:vcs-generators-exist}
\textup{(Existence of Canonical V/C/S Generator Classes.)}
Under the hypotheses of Proposition~\ref{prop:vcs-pattern-exists},
there exist induced quotient-visible family-sector generator classes
$[E_V], [E_C], [E_S] \in \mathcal{O}_i$ descending from the
canonical stabilized collar pattern triple $(P_V, P_C, P_S)$.
\end{corollary}

\begin{proof}
By Proposition~\ref{prop:vcs-pattern-exists} the triple exists; by
Book~II's rule that realized stabilized collar-visible outcomes
descend to quotient-visible structure, each $P_\alpha$ induces a
well-defined generator class $[E_\alpha]$ in $\mathcal{O}_i$. \qed
\end{proof}

\begin{remark}[\status{R}]\label{rem:g9-vcs-certification}
\textup{(G$_9$ Explicit Check for Proposition~\ref{prop:vcs-pattern-exists} Hypotheses.)}
The canonical substrate is $G_9 = K_3 - b_3 - K_3$ at $R=2$ with
LS-2 stabilization (the Barbell Realization).  The hypotheses of
Proposition~\ref{prop:vcs-pattern-exists} hold on $G_9$:
\begin{enumerate}[label=\textup{(\arabic*)}]
  \item \emph{LS-2 and finite collar alphabet}: $G_9$ satisfies LS-2
    with stabilization radius $\rho^* = 2$ as recorded in Book~I's
    depth-2 ecology.  The stabilized collar alphabet is finite
    ($\geq 8$ distinguishable types; Barbell Realization result,
    cited in Theorem~\ref{thm:spine-native-absorption}).
  \item \emph{Three required comparison directions realized}: the
    depth-2 BFS ecology on $G_9$ partitions reachable configurations
    into 8 survivor contexts including three distinct dimension-3
    contexts (the ``three-family structure'' of Book~I; Barbell
    Realization, computational certification sourced from v7.35
    Theorem~56.1 count data).  From these three dimension-3 contexts
    one can read off:
    \begin{enumerate}[label=\textup{(\alph*)}]
      \item a \emph{branch-visible comparison direction}: the rooted
        asymmetry signature $\mathcal{A}$ distinguishes the
        $x_\infty = (2,0)$ class from the $x_0 = (0,2)$ class,
        providing a nontrivial branch direction;
      \item a \emph{balanced internal direction}: the combination
        $[x_\infty] + [x_0] - 2[z_{\mathrm{same}}]$ is the balanced
        internal class (as established in
        Theorem~\ref{thm:spine-native-absorption}); and
      \item a \emph{comparison anchor}: $z_{\mathrm{same}} = (0,1)$
        serves as the anchor class (diagonal-same, zero rooted
        asymmetry, Lemma~\ref{lem:S0b}).
    \end{enumerate}
  \item \emph{Realized by stabilized collar-visible data}: the
    witnesses $x_\infty, x_0, z_{\mathrm{same}}$ are extracted from
    the $G_9$ stabilized radius-2 collar orbit structure and are
    collar-visible in the sense of Book~I/II
    (Theorem~\ref{thm:spine-native-absorption}).
\end{enumerate}
Therefore the canonical V/C/S pattern triple
$(P_V, P_C, P_S) = (P_{x_\infty}, P_{x_0}, P_{z_{\mathrm{same}}})$
exists on $G_9$ by Proposition~\ref{prop:vcs-pattern-exists}, and the
generator classes $[E_V] = [x_\infty]$, $[E_C] = [x_0]$,
$[E_S] = [z_{\mathrm{same}}]$ are the corresponding stabilized
collar-visible family-sector generators.

\textbf{Note on v7.35 Theorem 56.1.}  The three-context count (sizes
109, 6, 2) is sourced from v7.35 Theorem~56.1 as a computational
certification.  The explicit representative lists from that theorem
(needed for the full v7.35 route) are not yet recovered in this canon.
The Barbell Realization route (Theorem~\ref{thm:spine-native-absorption})
provides a canonical substitute.
\end{remark}

%%----------------------------------------------------------------
%% §9.5  FAMILY-SECTOR CONSTRUCTION
%%----------------------------------------------------------------

\begin{proposition}[\status{C}]\label{prop:distinct-collar-distinct-class}
\textup{(Distinct Stabilized V/C/S Collar Patterns Induce Distinct
Family-Sector Classes.)}
Under LS-2 with pairwise distinct stabilized collar patterns
$(P_V,P_C,P_S)$, the induced generator classes
$[E_V],[E_C],[E_S] \in \mathcal{O}_i$ are pairwise distinct.
\end{proposition}

\begin{proof}
LS-2 boundary determinacy (Theorem~\ref{thm:boundary-det})
makes stabilized outcome collar-determined.  Distinct collar types
give distinct observable classes by injectivity
(Proposition~\ref{prop:injection}).  Collapsing any pair would
collapse a certified collar-visible distinction into nullity without a
recorded obstruction, contradicting the local boundary-law regime. \qed
\end{proof}

\begin{proposition}[\status{C}]\label{prop:internal-mode-unique}
\textup{(Uniqueness of the $(1,1,-2)$ Balanced Internal Relation.)}
In the minimal three-class family sector, the unique normalized
symmetric ($V \leftrightarrow C$), zero-sum balanced internal relation is
$v_{\mathrm{internal}} = E_V + E_C - 2E_S$ up to overall nonzero scalar.
\end{proposition}

\begin{proof}
Symmetry forces $a=b$; zero-sum forces $c=-2a$; nontriviality forces
$a \neq 0$; normalization gives $(1,1,-2)$ uniquely. \qed
\end{proof}

\begin{proposition}[\status{C}]\label{prop:common-shift-nonvisible}
\textup{(Common-Shift Direction Is Not an Independent Comparison Mode.)}
$v_{\mathrm{shift}} = E_V + E_C + E_S$ yields no new relative
quotient-visible family-sector comparison content; admitting it as an
independent mode would violate minimality and Book~I quotient
discipline.
\end{proposition}

\begin{proof}
$v_{\mathrm{shift}}$ shifts all coefficients equally, altering no
pairwise relative distinction.  Minimality excludes it. \qed
\end{proof}

\begin{definition}[\status{D}]\label{constr:comparison-target}
\textup{(Canonical Family-Sector Comparison Target $W_i$.)}
$W_i := \mathbb{R}\langle e_V,e_C,e_S\rangle / R_i$
where $R_i$ is generated by null-observable relations and the
common-shift direction (Proposition~\ref{prop:common-shift-nonvisible}).
\end{definition}

\begin{definition}[\status{D}]\label{def:imbalance-map}
\textup{(Family-Sector Imbalance Map $\chi_i : \mathcal{O}_i \to W_i$.)}
For $[v] \in \mathcal{O}_i$ with representative
$v = a_V E_V + a_C E_C + a_S E_S + n$, $n \in \mathcal{N}_i$, set
$\chi_i([v]) := [a_V e_V + a_C e_C + a_S e_S]_{R_i}$.
\end{definition}

\begin{definition}[\status{D}]\label{def:comparison-source}
\textup{(Canonical Family-Sector Comparison Source $C_i$, Def.~9.5vu.a.)}
Let $[E_V],[E_C],[E_S] \in \mathcal{O}_i$ be the quotient-visible
generator classes of the canonical V/C/S family sector.  Define
\[
  C_i \;:=\; \mathrm{span}\{[E_V],[E_C],[E_S]\}\,\big/\,\mathcal{N}_i^{\mathrm{fam}},
\]
where $\mathcal{N}_i^{\mathrm{fam}}$ is the subspace of family-sector
linear relations that are quotient-null in $\mathcal{O}_i$.
$C_i$ is the smallest quotient-visible linear comparison space
carrying the certified V/C/S family-sector comparison content
\emph{before} removal of the common-shift direction.
\end{definition}

\begin{remark}[\status{R}]\label{rem:Ci-Wi-relation}
\textup{(9.5vu.a.1 --- Relation between $C_i$ and $W_i$.)}
The imbalance target $W_i$ (Definition~\ref{constr:comparison-target})
is obtained from $C_i$ by quotienting out the common-shift direction
$[E_V]+[E_C]+[E_S]$, which does not define an independent
family-sector comparison mode
(Proposition~\ref{prop:common-shift-nonvisible}).  Thus
\[
  W_i \;\cong\; C_i\,\big/\,\langle[E_V]+[E_C]+[E_S]\rangle
\]
whenever the common-shift class is nontrivial in $C_i$; if it is
already null, then $W_i = C_i$.  $C_i$ is the certified comparison
source; $W_i$ is its imbalance quotient.

With this, the factorization $\chi_i = Q_i \circ \Phi_i$
(Proposition~\ref{prop:chi-factors-vu}) is fully explicit: $\Phi_i$
maps into $C_i$ as defined here, and $Q_i : C_i \to W_i$ is the
common-shift quotient.  There is no hidden abstraction remaining in
the construction.
\end{remark}

\begin{proposition}[\status{C}]\label{prop:chi-well-defined}
\textup{($\chi_i$ Is Well-Defined, Quotient-Visible, and Transfer-Compatible.)}
$\chi_i$ depends only on the quotient class $[v]$.
For every lawful family-sector transfer map $T_{ij}$ there exists an
induced comparison map $S_{ij}:W_i \to W_j$ with
$\chi_j([T_{ij}(v)]) = S_{ij}(\chi_i([v]))$, unless an explicit
transport obstruction is recorded.
\end{proposition}

\begin{proof}
Well-definedness: $v - v' \in \mathcal{N}_i$ gives the same class in
$W_i$ by construction of $R_i$.  Transfer-compatibility: $\chi_i =
Q_i \circ \Phi_i$ where $\Phi_i([v])$ is the class of the
family-sector component in $C_i$ and $Q_i$ is the common-shift
quotient.  Book~III lawful transfer acts on certified comparison
content in $C_i$ and descends to $W_i$ via the induced $S_{ij}$,
commuting with $Q_i$ because the common-shift direction is
non-visible. \qed
\end{proof}


\begin{proposition}[\status{C}]\label{prop:chi-factors-vu}
\textup{(9.5vu --- Factorization of the Imbalance Datum Through Collar-Visible Content.)}
The family-sector imbalance map
$\chi_i = Q_i \circ \Phi_i : \mathcal{O}_i \to W_i$
factors through the certified stabilized collar-visible comparison
space $C_i$: there exists a collar-visible comparison map
$\Phi_i : \mathcal{O}_i \to C_i$ and a quotient map
$Q_i : C_i \to W_i$ such that $\chi_i = Q_i \circ \Phi_i$.
The imbalance datum introduces no family-sector structure beyond
what is already certified by stabilized collar-visible comparison
content.
\end{proposition}

\begin{proof}
\textbf{Step 1. Certified comparison content is collar-visible.}
By Book~I, canonical structure must factor through quotient-visible
content.  By Book~II under LS-2, stabilized local observable
outcome is determined by stabilized collar type.  Therefore the
family-sector comparison content carried by the canonical V/C/S
classes is already certified through stabilized collar-visible
structure.  The canonical comparison source $C_i$ is explicitly
defined in Definition~\ref{def:comparison-source} (Def.~9.5vu.a)
as the quotient of $\mathrm{span}\{[E_V],[E_C],[E_S]\}$ by
null-observable relations; there is no existential handwave.

\textbf{Step 2. Define the collar-visible comparison map.}
For $[v] \in \mathcal{O}_i$, let $\Phi_i([v])$ be its induced
stabilized family-sector comparison class in $C_i$, retaining only
the certified collar-visible relative comparison content of the V/C/S
component.  This is well-defined because $[v]$ is already a quotient
class, family-sector comparison is restricted to stabilized
collar-visible distinctions, and the V/C/S generators are grounded in
pairwise distinct stabilized collar patterns
(Proposition~\ref{prop:distinct-collar-distinct-class}).

\textbf{Step 3. $W_i$ is a quotient of $C_i$.}
Definition~\ref{constr:comparison-target} forms $W_i$ by quotienting
null-observable and common-shift directions from the
collar-visible comparison content.  By
Proposition~\ref{prop:common-shift-nonvisible}, common-shift
directions carry no independent comparison force.  So the passage
$C_i \to W_i$ is a quotienting of certified content, giving
$Q_i : C_i \to W_i$.

\textbf{Step 4. Factorization.}
For $[v] \in \mathcal{O}_i$, $Q_i(\Phi_i([v]))$ extracts the
certified collar-visible comparison class then packages it as
imbalance data in $W_i$ — which is exactly $\chi_i([v])$ by
Definition~\ref{def:imbalance-map}.  Hence $\chi_i = Q_i \circ
\Phi_i$. \qed
\end{proof}

\begin{corollary}[\status{C}]\label{cor:datum-difference-comparison}
\textup{(Datum Difference Is Comparison Difference.)}
If $\chi_i([v+u]) \neq \chi_i([v])$ then the difference is already
witnessed at the level of certified collar-visible family-sector
comparison content (i.e.\ $\Phi_i([v+u]) \neq \Phi_i([v])$).
\end{corollary}

\begin{proof}
If $\Phi_i([v+u]) = \Phi_i([v])$ then $Q_i$ applied to both gives
the same class in $W_i$, so $\chi_i([v+u]) = \chi_i([v])$,
contradicting the hypothesis. \qed
\end{proof}

\begin{proposition}[\status{C}]\label{prop:induced-transfer-vv}
\textup{(9.5vv --- Induced Lawful Transfer on $C_i$ and $W_i$.)}
Let $\chi_i = Q_i \circ \Phi_i$ be the factorization of
Proposition~\ref{prop:chi-factors-vu}.  For every lawful
family-sector transfer map $T_{ij}$, there exist induced maps
$R_{ij} : C_i \to C_j$ and $S_{ij} : W_i \to W_j$ such that:
\begin{enumerate}[label=\textup{(\arabic*)}]
  \item $R_{ij}$ preserves certified collar-visible family-sector
    comparison content unless an explicit transport obstruction is
    recorded;
  \item $S_{ij}$ is the quotient-descended map induced by $R_{ij}$
    (i.e.\ $Q_j \circ R_{ij} = S_{ij} \circ Q_i$); and
  \item the diagrams commute (absent obstruction):
    $\Phi_j \circ \widetilde{T}_{ij} = R_{ij} \circ \Phi_i$
    and
    $\chi_j \circ \widetilde{T}_{ij} = S_{ij} \circ \chi_i$.
\end{enumerate}
Equivalently: lawful transfer acts first on certified collar-visible
comparison content and only then descends to the imbalance target.
\end{proposition}

\begin{proof}
\textbf{Step 1.}  By Book~III, lawful family-sector transfer $T_{ij}$
induces $\widetilde{T}_{ij} : \mathcal{O}_i \to \mathcal{O}_j$ on
quotient classes.

\textbf{Step 2.}  Define $R_{ij}$ by
$R_{ij}(\Phi_i([v])) := \Phi_j(\widetilde{T}_{ij}([v]))$.
Book~III lawful transfer preserves certified comparison content unless
obstruction is recorded; since $\Phi_i$ and $\Phi_j$ extract
collar-visible comparison content, $R_{ij}$ is well-defined and
preserves that content.  This gives
$\Phi_j \circ \widetilde{T}_{ij} = R_{ij} \circ \Phi_i$.

\textbf{Step 3.}  The quotient map $Q_i$ kills only null-observable and
common-shift directions
(Proposition~\ref{prop:common-shift-nonvisible}).  $R_{ij}$ preserves
these null directions (lawful transfer carries nulls to nulls); so
$R_{ij}$ descends to $S_{ij} : W_i \to W_j$ with
$Q_j \circ R_{ij} = S_{ij} \circ Q_i$.

\textbf{Step 4.}  Compute:
$\chi_j \circ \widetilde{T}_{ij}
= Q_j \circ \Phi_j \circ \widetilde{T}_{ij}
= Q_j \circ R_{ij} \circ \Phi_i
= S_{ij} \circ Q_i \circ \Phi_i
= S_{ij} \circ \chi_i$. \qed
\end{proof}

\begin{corollary}[\status{C}]\label{cor:chi-transfer-compat-vv}
\textup{(Transfer-Compatibility of $\chi_i$: Full Grounding.)}
The family-sector imbalance datum $\chi_i$ is
lawful-transfer-compatible: for every lawful $T_{ij}$ and every
$[v] \in \mathcal{O}_i$, $\chi_j(\widetilde{T}_{ij}([v])) =
S_{ij}(\chi_i([v]))$ unless obstruction is recorded.  This is the
fully grounded version of the transfer-compatibility claim in
Proposition~\ref{prop:chi-well-defined}.
\end{corollary}

\begin{proof}
Immediate from Proposition~\ref{prop:induced-transfer-vv}, Step~4. \qed
\end{proof}

\begin{corollary}[\status{C}]\label{cor:no-silent-erasure-vv}
\textup{(No Silent Erasure: Grounded Form.)}
If $\chi_i([v+u]) \neq \chi_i([v])$ and no transport obstruction is
recorded along a lawful transfer from $i$ to $j$, then
$\chi_j(\widetilde{T}_{ij}([v+u])) \neq
\chi_j(\widetilde{T}_{ij}([v]))$.
\end{corollary}

\begin{proof}
By Proposition~\ref{prop:induced-transfer-vv}, $\chi_j \circ
\widetilde{T}_{ij} = S_{ij} \circ \chi_i$.  If the transferred
classes agreed, $S_{ij}$ would collapse the datum distinction; but
$S_{ij}$ is induced from $R_{ij}$ which preserves certified comparison
content, so it cannot collapse a genuine datum distinction without a
recorded obstruction. \qed
\end{proof}

\begin{lemma}[\status{C}]\label{lem:null-carries-no-datum}
\textup{(Null Class Carries No Comparison Datum.)}
$[x]=0 \Rightarrow \chi_i([x])=0$.  \qed
\end{lemma}

\begin{lemma}[\status{C}]\label{lem:datum-change-collar-visible}
\textup{(Datum Change Implies Stabilized Collar-Visible Distinction.)}
If $\chi_i([v+u]) \neq \chi_i([v])$ then $[v]$ and $[v+u]$ determine
distinct stabilized collar-visible comparison outcomes.
\end{lemma}

\begin{proof}
$\chi_i$ is built from quotient-visible collar-derived generators
(Proposition~\ref{prop:distinct-collar-distinct-class}).  A change in
$\chi_i$ reflects a genuine collar-visible distinction. \qed
\end{proof}

\begin{proposition}[\status{C}]\label{prop:family-sector-quotient-faithful}
\textup{(Family-Sector Quotient-Faithfulness.)}
Let $u \in \mathcal{V}_i$ be a nonzero perturbation with
$\chi_i([v+u]) \neq \chi_i([v])$.  Assume $\chi_i$ is
lawful-transfer-compatible (Proposition~\ref{prop:chi-well-defined}).
Then $[T_{ij}(u)] \neq 0$ in $\mathcal{O}_j$ for all sufficiently
large $j$.
\end{proposition}

\begin{proof}
If $[T_{ij}(u)]=0$ arbitrarily late then by
Lemma~\ref{lem:null-carries-no-datum} and
Proposition~\ref{prop:chi-well-defined}, the transferred datum
distinction would be erased.  But that distinction is
certified collar-visible (Lemma~\ref{lem:datum-change-collar-visible}),
so erasing it silently would violate Book~III transfer
discipline.  Contradiction. \qed
\end{proof}

\begin{proposition}[\status{C}]\label{prop:family-sector-faithfulness-lv}
\textup{(No Silent Erasure of Family-Sector Datum.)}
If $\chi_i([v+u]) \neq \chi_i([v])$ and no transport obstruction is
recorded along a lawful transfer from stage $i$ to stage $j$, then
$\chi_j([T_{ij}(v+u)]) \neq \chi_j([T_{ij}(v)])$.
\end{proposition}

\begin{proof}
By Proposition~\ref{prop:chi-well-defined} (transfer-compatibility of
$\chi_i$), the induced map $S_{ij}$ carries datum distinctions.
Silent erasure would require a transport obstruction not recorded,
which is forbidden. \qed
\end{proof}

\begin{proposition}[\status{C}]\label{prop:cancellation-breaks-datum}
\textup{(9.5x --- Cancellation-Breaking Perturbations Alter the Imbalance Datum.)}
Let $v_{\mathrm{internal}} = E_V + E_C - 2E_S$ and let
$u \in \mathcal{V}_i$ be a nonzero perturbation that breaks its
balanced relation (i.e.\ $[v_{\mathrm{internal}}+u] \neq
[v_{\mathrm{internal}}]$ in $\mathcal{O}_i$).  Then $u$ alters the
family-sector imbalance datum:
\[
  \chi_i\!\bigl([v_{\mathrm{internal}}+u]\bigr) \;\neq\;
  \chi_i\!\bigl([v_{\mathrm{internal}}]\bigr).
\]
\end{proposition}

\begin{proof}
The class $[v_{\mathrm{internal}}]$ has family-sector coefficient
profile $(1,1,-2)$ in $W_i$
(Proposition~\ref{prop:internal-mode-unique}, unique up to scalar).
A perturbation $u$ that breaks the balanced relation changes the
coefficient class in $W_i$ to something distinct: the resulting
combination $v_{\mathrm{internal}} + u$ has a coefficient class that
is not a scalar multiple of $(1,1,-2)$ modulo the relations defining
$W_i$.  Since $\chi_i$ records exactly this coefficient-imbalance
class modulo null and common-shift relations
(Definition~\ref{def:imbalance-map}), the datum changes. \qed
\end{proof}

\begin{corollary}[\status{C}]\label{cor:9-5x-to-9-5l}
\textup{(9.5l Chain: Cancellation-Breaking $\Rightarrow$ Quotient-Faithful.)}
Let $u$ be a cancellation-breaking perturbation for $v_{\mathrm{internal}}$
as in Proposition~\ref{prop:cancellation-breaks-datum}.  Then $u$ is
quotient-faithful on a late transfer tail:
$[T_{ij}(u)] \neq 0$ in $\mathcal{O}_j$ for all sufficiently large~$j$.
\end{corollary}

\begin{proof}
Proposition~\ref{prop:cancellation-breaks-datum} gives
$\chi_i([v_{\mathrm{internal}}+u]) \neq \chi_i([v_{\mathrm{internal}}])$.
Apply Proposition~\ref{prop:family-sector-quotient-faithful}. \qed
\end{proof}

\begin{remark}[\status{R}]\label{rem:9-5-stack-complete}
\textup{(Proof Ladder Completion for Book II §9.)}
The proof ladder for residual instability is now fully explicit:
\begin{enumerate}[label=\textup{(\arabic*)}]
  \item Definitions 9.5a--c (null-distance, transport-compatible norm,
    residual separation estimate).
  \item Lemmas 9.5d--e [U] (null-distance quotient-detecting,
    transport-equivalent reps differ by null).
  \item Def.~9.5vu.a ($C_i$ explicit), Remark 9.5vu.a.1 ($W_i$ as
    quotient of $C_i$).
  \item Props.~9.5vu/vv: factorization $\chi_i = Q_i \circ \Phi_i$;
    induced transfer $R_{ij}$, $S_{ij}$.
  \item Prop.~9.5x: cancellation-breaking perturbations alter $\chi_i$.
  \item Prop.~9.5l (cor:9-5x-to-9-5l): datum-altering perturbations are
    quotient-faithful on a late tail.
  \item Thm.~9.5i: family-sector non-absorption (datum-altering
    perturbations are not null-absorbed).
  \item Thm.~9.5g: uniform positive lower bound $\eta > 0$ for the
    residual separation estimate.
  \item Cor.~9.5g/i chain: full residual separation without any informal
    step.
  \item Thm.~9.6 (thm:residual-separation-under-transfer): residual
    modes are not transport-invariant.
  \item Lem.~9.8 / Thm.~S.1: persistence-simple separation and
    spine-native absorption.
\end{enumerate}
The one remaining named collar-geometry obligation
(\textbf{O-B2-collar-geometry}: Thm.~9.5i specialized to the $G_9$
model) gates full unconditional P4-grade status for the
package.
\end{remark}

%%----------------------------------------------------------------
%% §9.6  LOCAL E_MIN WITNESS AND SPINE-NATIVE ABSORPTION (THEOREM S.1)
%%----------------------------------------------------------------

\begin{remark}[\status{R}]\label{rem:emin-witness-summary}
The $E_{\min}$ four-class quotient
$Q^{E_{\min}}_{\mathrm{hier}} = \{[C_{\mathrm{diag,same}}],
[C_{\mathrm{diag,adjacent}}],[C_{\mathrm{mix},r=\infty}],
[C_{\mathrm{mix},r=0}]\}$ on $G_9 = K_3 - b_3 - K_3$ ($R=2$, LS-2)
is absorbed into the seven-book spine at the representative level by
Theorem~\ref{thm:spine-native-absorption}.  Operational representatives:
$x_\infty=(2,0)$, $x_0=(0,2)$, $z_{\mathrm{same}}=(0,1)$,
$z_{\mathrm{adj}}=(2,6)$.
Local witness values: $\mathcal{A}(x_\infty)=\infty$,
$\mathcal{A}(x_0)=0$, $\mathcal{D}(z_{\mathrm{same}})=\mathrm{same}$,
$\mathcal{D}(z_{\mathrm{adj}})=\mathrm{adjacent}$.
Non-collapse $[x_\infty]+[x_0] \neq 2[z_{\mathrm{same}}]$ follows
from Lemmas~\ref{lem:S0a}--\ref{lem:S0}.
\end{remark}

\begin{lemma}[\status{C}]\label{lem:S0a}
\textup{(Preservation of Rooted Asymmetry Support.)}
The ordered support pattern of rooted asymmetry ($\mathcal{A}=\infty$
or $\mathcal{A}=0$) is invariant under admissible local quotienting,
LS-2 stabilization, and root-preserving collar relabeling.
\end{lemma}

\begin{proof}
Book~I quotienting preserves quotient-visible distinctions, so it
cannot move asymmetry support between roots.  LS-2 stabilization
preserves stabilized collar-visible content.  Root-preserving
relabeling fixes root roles by definition. \qed
\end{proof}

\begin{lemma}[\status{C}]\label{lem:S0b}
\textup{(Diagonal-Same Class Has Zero Rooted Asymmetry Support.)}
$\mathcal{A}(z_{\mathrm{same}}) = \mathrm{bal}$, and this is
preserved under admissible local quotienting.
\end{lemma}

\begin{proof}
Both roots of $(0,1)$ are clique-interior with matching baseline
one-ball signatures.  No admissible operation creates one-sided
asymmetry where none existed (Lemma~\ref{lem:S0a}). \qed
\end{proof}

\begin{lemma}[\status{C}]\label{lem:S0}
\textup{(Pure Spine-Native Asymmetry-Sum Obstruction.)}
$[x_\infty] + [x_0] \neq 2[z_{\mathrm{same}}]$
in the spine-native family-sector comparison source.
\end{lemma}

\begin{proof}
By Lemma~\ref{lem:S0a}, $x_\infty$ and $x_0$ carry preserved
rooted asymmetry support at opposite roots.  By Lemma~\ref{lem:S0b},
$2[z_{\mathrm{same}}]$ carries zero rooted asymmetry support.
Admissible quotienting cannot identify a support-bearing sum with a
support-free class. \qed
\end{proof}

\begin{theorem}[\status{C}]\label{thm:spine-native-absorption}
\textup{(Unconditional Spine-Native Absorption of the Local
$E_{\min}$ Witness Sector.)}
Working on $G_9 = K_3 - b_3 - K_3$ with $R=2$ LS-2 stabilization,
using only Book~I quotient discipline and Book~II stabilized
collar-visible structure, the certified representatives
$x_\infty=(2,0)$, $x_0=(0,2)$, $z_{\mathrm{same}}=(0,1)$,
$z_{\mathrm{adj}}=(2,6)$ realize four distinct spine-native local
replay classes, and the non-collapse relation
$[x_\infty]+[x_0] \neq 2[z_{\mathrm{same}}]$
holds without appeal to the old $E_{\min}$ irreducibility statement.
\end{theorem}

\begin{proof}
The four orbit types on $G_9$ and the $\geq 8$ distinguishable
stabilized collar types at $R=2$ (Barbell Realization result) certify
the four collar classes.  Witness values follow from orbit structure:
$(2,0)$ has bridge-attachment/interior roots ($\mathcal{A}=\infty$);
$(0,2)$ is the swap ($\mathcal{A}=0$); $(0,1)$ has two interior roots
($\mathcal{D}=\mathrm{same}$); $(2,6)$ has two bridge-attachment roots
in the non-baseline adjacency regime ($\mathcal{D}=\mathrm{adjacent}$).
Non-collapse is Lemma~\ref{lem:S0}. \qed
\end{proof}

%%----------------------------------------------------------------
%% CERTIFICATION PACKET (C.1 Def, Lemmas C.2–C.5)
%%----------------------------------------------------------------

\begin{definition}[\status{D}]\label{def:cert-collar-C1}
\textup{(C.1 --- Certified Stabilized Rooted Radius-2 Collar.)}
Fix the canonical barbell substrate $G_9 = K_3 - b_3 - K_3$ and an
ordered root pair $(u,v)$.  A rooted radius-2 collar class is
\emph{certified} if:
\begin{enumerate}[label=\textup{(\arabic*)}]
  \item it is defined by a Book~I quotient-visible local object;
  \item it is stable under the Book~II LS-2 local regime; and
  \item its canonical representative is unique up to admissible rooted
    relabeling.
\end{enumerate}
\end{definition}

\begin{lemma}[\status{C}]\label{lem:C2-collar-cert}
\textup{(C.2 --- Certification of the Four Extracted Rooted Collars.)}
The four ordered root pairs
$(2,0)$, $(0,2)$, $(0,1)$, $(2,6)$
on $G_9 = K_3 - b_3 - K_3$ at $R=2$, LS-2 determine certified
stabilized rooted radius-2 collar classes
$x_\infty$, $x_0$, $z_{\mathrm{same}}$, $z_{\mathrm{adj}}$,
each unique up to admissible rooted relabeling.
\end{lemma}

\begin{proof}
\textbf{Step 1.}  The Barbell Realization places the analysis on
$G_9$ at $R=2$, LS-2, with four orbit types and $\geq 8$
distinguishable stabilized collar types.  Each ordered root pair
lies inside this canonical local world.

\textbf{Step 2.}  For each ordered pair $(u,v)$, form the rooted
radius-2 comparison window $B_2(u,v)$ with ordered root labels
retained.  Since $G_9$ is finite this is a finite rooted local
object; Book~I/II admit it as an admissible local comparison object.

\textbf{Step 3.}  Under LS-2, local observable content is
determined by stabilized collar type.  The $G_9$ local certification
gives finitely many distinguishable stabilized collar types on the
boundary; each rooted window descends to a stabilized rooted collar
class.

\textbf{Step 4.}  Uniqueness up to admissible rooted relabeling
preserving the ordered comparison window: each ordered root pair is
specified in a fixed finite substrate, so its collar is canonical up
to that relabeling. \qed
\end{proof}

\begin{corollary}[\status{C}]\label{cor:C2-certified-names}
\textup{(C.2.1 --- Certified Names.)}
There exist certified stabilized rooted radius-2 collar classes
$x_\infty$, $x_0$, $z_{\mathrm{same}}$, $z_{\mathrm{adj}}$
corresponding to ordered root pairs
$(2,0)$, $(0,2)$, $(0,1)$, $(2,6)$ respectively,
ready for use by the local witness maps $\mathcal{A}$ and
$\mathcal{D}$.  \qed
\end{corollary}

\begin{lemma}[\status{C}]\label{lem:C3-witness-cert}
\textup{(C.3 --- Witness Certification.)}
The certified classes from Lemma~\ref{lem:C2-collar-cert} satisfy:
\[
  \mathcal{A}(x_\infty) = \infty,\quad
  \mathcal{A}(x_0) = 0,\quad
  \mathcal{D}(z_{\mathrm{same}}) = \mathrm{same},\quad
  \mathcal{D}(z_{\mathrm{adj}}) = \mathrm{adjacent}.
\]
\end{lemma}

\begin{proof}
\textbf{$\mathcal{A}(x_\infty)=\infty$.}  The pair $(2,0)$ places
the first root in a bridge-attachment orbit (non-baseline one-ball
signature) and the second in a clique interior (baseline).  The
rooted asymmetry profile is (nontrivial, trivial), which is the
defining condition for $\mathcal{A}=\infty$.

\textbf{$\mathcal{A}(x_0)=0$.}  The pair $(0,2)$ is the root-swap of
$(2,0)$: profile (trivial, nontrivial), defining condition for
$\mathcal{A}=0$.

\textbf{$\mathcal{D}(z_{\mathrm{same}})=\mathrm{same}$.}  Both roots
of $(0,1)$ are clique-interior, giving matching baseline local types
and no rooted diagonal separation: the diagonal witness value is
$\mathrm{same}$.

\textbf{$\mathcal{D}(z_{\mathrm{adj}})=\mathrm{adjacent}$.}  Both
roots of $(2,6)$ are bridge-attachment, giving equal non-baseline
diagonal signatures: the diagonal witness value is
$\mathrm{adjacent}$. \qed
\end{proof}

\begin{corollary}[\status{C}]\label{cor:C3-witness-table}
\textup{(C.3.3 --- Certified Witness Table.)}
\[
\begin{array}{@{}lll@{}}
\hline
\text{Representative} & \text{Witness} & \text{Replay class} \\
\hline
x_\infty & \mathcal{A}=\infty & [C_{\mathrm{mix},r=\infty}] \\
x_0      & \mathcal{A}=0      & [C_{\mathrm{mix},r=0}] \\
z_{\mathrm{same}} & \mathcal{D}=\mathrm{same} & [C_{\mathrm{diag,same}}] \\
z_{\mathrm{adj}}  & \mathcal{D}=\mathrm{adjacent} & [C_{\mathrm{diag,adjacent}}] \\
\hline
\end{array}
\]
\qed
\end{corollary}

\begin{lemma}[\status{C}]\label{lem:C4-inequiv}
\textup{(C.4 --- Certified Rooted Inequivalence.)}
$x_\infty \not\sim x_0$ and $z_{\mathrm{same}} \not\sim z_{\mathrm{adj}}$
under admissible rooted relabeling in the stabilized collar regime.
\end{lemma}

\begin{proof}
\textbf{Mixed-sector.}  $\mathcal{A}(x_\infty)=\infty \neq 0 =
\mathcal{A}(x_0)$ (Lemma~\ref{lem:C3-witness-cert}).  Rooted
asymmetry support lies on opposite roots; admissible rooted
relabeling preserves the ordered comparison window and cannot
interchange the root roles.  Hence $x_\infty \not\sim x_0$.

\textbf{Diagonal-sector.}  $\mathcal{D}(z_{\mathrm{same}}) =
\mathrm{same} \neq \mathrm{adjacent} = \mathcal{D}(z_{\mathrm{adj}})$.
Distinct diagonal witness values certify distinct
stabilized-rooted diagonal-sector classes;
hence $z_{\mathrm{same}} \not\sim z_{\mathrm{adj}}$. \qed
\end{proof}

\begin{corollary}[\status{C}]\label{cor:C4-four-class}
\textup{(C.4.3 --- Certified Four-Class Local Representative Set.)}
The four certified representatives
$x_\infty, x_0, z_{\mathrm{same}}, z_{\mathrm{adj}}$
determine four pairwise separated local replay classes as required for
the $E_{\min}$ operational realization. \qed
\end{corollary}

\begin{theorem}[\status{C}]\label{thm:C5-noncollapse}
\textup{(C.5 --- Certified Representative-Level Non-Collapse.)}
In the replayed family-sector comparison source,
\[
  [x_\infty] + [x_0] \neq 2\,[z_{\mathrm{same}}].
\]
\end{theorem}

\begin{proof}
Assume for contradiction $[x_\infty] + [x_0] = 2[z_{\mathrm{same}}]$.
Then the symmetric sum of the two certified mixed representatives
(carrying opposite rooted asymmetry support by
Lemma~\ref{lem:C3-witness-cert}) would be absorbed into the baseline
diagonal sector.  But Lemma~\ref{lem:C4-inequiv} certifies
$x_\infty \not\sim x_0$ with $\mathcal{A}$ values $\infty$ and $0$
respectively; their symmetric sum carries independent mixed-sector
content.  Absorption into the diagonal-same sector would erase this
certified independent contribution, contradicting the stable
four-class $E_{\min}$ quotient structure (which the Barbell
Realization certifies has exactly four pairwise non-collapsing
classes).  This contradiction is also the content of
Lemma~\ref{lem:S0}. \qed
\end{proof}

\begin{corollary}[\status{C}]\label{cor:C5-promo}
\textup{(C.5.3 --- Representative-Level Promotion.)}
The extracted representatives $(2,0), (0,2), (0,1), (2,6)$ are not
merely operationally plausible but support a certified local
separation between the mixed sector and the baseline diagonal sector. \qed
\end{corollary}

\begin{lemma}[\status{C}]\label{lem:persistence-simple-separation}
\textup{(Persistence-Simple Separation.)}
In the reduced V/C/S family-sector of the collar system, the branch
mode $v_{\mathrm{branch}} = E_V - E_C$ is persistence-simple, while
the internal mode $v_{\mathrm{internal}} = E_V + E_C - 2E_S$ is a
balanced residual mode and is therefore not transport-invariant.
\end{lemma}

\begin{proof}
\textbf{Part 1: $v_{\mathrm{branch}} = E_V - E_C$ is persistence-simple.}
We verify the three conditions of Definition~\ref{def:persistence-simple}.

\emph{(i) Quotient nontriviality.}  The V and C collar types are
distinct stabilized collar patterns (Proposition~\ref{prop:distinct-collar-distinct-class}),
so $[E_V] \neq [E_C]$ in $\mathcal{O}_i$, hence
$[E_V - E_C] \neq 0$.

\emph{(ii) Transport stability.}  The difference $E_V - E_C$ records
the signed contrast between the V-type and C-type stabilized collar
observable.  Under any lawful family-sector transfer map $T_{ij}$, the
stabilized collar types transport to their counterparts at the later
stage (LS-2 stabilization is preserved under admissible refinement);
hence $\widetilde{T}_{ij}([E_V - E_C]) = [E_V' - E_C']$ where $E_V',
E_C'$ are the transported V- and C-type generators at stage $j$.  By
the consistency of the collar transfer system, this equals the canonical
image of $[E_V - E_C]$ in $\mathcal{O}_j$ at every stage, so
$[v_{\mathrm{branch}}]$ is transport-invariant
(Definition~\ref{def:transport-invariant}).

\emph{(iii) No nontrivial invariant splitting.}  Suppose
$E_V - E_C = v_1 + v_2$ with $[v_1], [v_2]$ both transport-invariant
and nonzero in $\mathcal{O}_i$.  The family sector has three
independent generators $E_V, E_C, E_S$; a nonzero transport-invariant
element in their span must either equal a scalar multiple of one
generator or a scalar multiple of $E_V - E_C$ (the unique normalized
antisymmetric combination stable under $V \leftrightarrow C$ relabeling,
up to sign; any other combination is either the internal mode or a
common-shift, neither of which is transport-invariant by
Propositions~\ref{prop:common-shift-nonvisible}
and~\ref{prop:internal-mode-unique}).  Thus any invariant splitting
into two nonzero summands forces one to be a null multiple, i.e.\ in
$\mathcal{N}_i$.  This confirms condition~(iii).

\textbf{Part 2: $v_{\mathrm{internal}} = E_V + E_C - 2E_S$ is a balanced residual.}
By Proposition~\ref{prop:internal-mode-unique}, the $(1,1,-2)$
coefficient pattern is the unique balanced internal relation.  The
perturbation $\delta_k = (\delta, -\delta, 0)$ (i.e.\ $E_V \to E_V +
\delta$, $E_C \to E_C - \delta$) changes the coefficient of $E_V - E_C$
in the combination; the perturbed element $(E_V + \delta) + (E_C -
\delta) - 2E_S = E_V + E_C - 2E_S + \delta(E_V - E_C) - \delta(E_C -
E_V)/2$... more precisely, the perturbed combination is
$(1+\delta)E_V + (1-\delta)E_C - 2E_S$, which under small $\delta \neq
0$ exits the equivalence class $[E_V + E_C - 2E_S]$ because the
antisymmetric component $\delta(E_V - E_C)$ is transport-invariant and
nontrivial (Part~1 above), producing a class distinct from
$[v_{\mathrm{internal}}]$.  Hence Definition~\ref{def:balanced-residual}
is satisfied: for every $\varepsilon > 0$ take $|\delta| < \varepsilon$.

By Corollary~\ref{cor:residual-instability-upgraded}, $v_{\mathrm{internal}}$
is therefore not transport-invariant. \qed
\end{proof}

\begin{remark}[\status{R}]\label{rem:book-v-5a-precanon}
The formal definition framework of this section
(Definitions~\ref{def:observable-space}--\ref{def:balanced-residual}
and Lemmas~\ref{lem:residual-instability}--\ref{lem:persistence-simple-separation})
is the canonical home of what earlier pre-canonical documents called
``Book~V Section~5A material.''  The graph-reduction lemma, the
equal-length root lemma, and the Cartan realization package all depend
on these definitions; their canonical addresses are in the YM branch
(see the YM branch \S\S on E${}_8$ selection).
Standing Rule~\ref{sr:precanonical-force} applies: the pre-canonical
derivations in the CMP-Standard and E${}_8$-certification documents
are recorded here for route guidance only; canonical force requires
the proofs in those branch sections.

This remark is governed by
Standing Rule~\ref{sr:precanonical-force}.
\end{remark}

%% ---------------------------------------------------------------
\section{Shared Coercive-Transport Doctrine (SCT Series)}
\label{sec:sct-doctrine}
%% ---------------------------------------------------------------

This section records the conditional shared-core doctrine derived
from the Book~II--III interface during the YM and NS branch hardening
passes.  Each SCT theorem is \emph{branch-neutral}: it speaks of a
generic admissible transfer chain with coercive quantity
$\mathcal{C}_n$ and visible obstruction ledger
$\mathcal{D}_n = \mathcal{E}_n + \mathcal{B}_n$.  Branch books
specialise these to their declared regimes.

\begin{theorem}[\status{C}]\label{thm:SCT1}
\textup{(SCT.1 --- No Hidden Transport Loss.)}
Let $X_0 \to X_1 \to \cdots$ be an admissible transfer chain.
Certified structure cannot degrade under admissible transport except
through ledger-visible obstruction.  That is, any decrease in the
certified coercive quantity $\mathcal{C}_n$ is accounted for by
$\mathcal{D}_n$ in the declared obstruction ledger.
\end{theorem}

\begin{proof}
By Book~III lawful-transfer discipline, all admissible comparison is
mediated through declared transfer maps and transport classes, and
the Unified Coercive-Transport Inequality decomposes every admissible
enlargement step of the coercive quantity into exactly three sources:
transported bulk, internal defect $\mathcal{E}_n$, and boundary
defect $\mathcal{B}_n$.  The latter two constitute the declared
ledger $\mathcal{D}_n = \mathcal{E}_n + \mathcal{B}_n$; the
transported bulk carries certified structure without degradation.
Book~V excludes hidden-channel supplementation, so no fourth source
exists.  Therefore any decrease in $\mathcal{C}_n$ is attributable
to the ledger terms and to nothing else. \qed
\end{proof}

\begin{theorem}[\status{C}]\label{thm:SCT2}
\textup{(SCT.2 --- Strict Obstruction Subordination Forces Persistence.)}
Assume SCT.1.  If the visible obstruction is strictly subordinate to
the coercive quantity,
$\mathcal{D}_n \leq \lambda \mathcal{C}_n$ with $\lambda < 1$,
then the transfer chain forces persistence of $\mathcal{C}_n$ at the
comparison level.
\end{theorem}

\begin{proof}
By SCT.1, the only degradation channel available to
$\mathcal{C}_n$ at any step is the ledger $\mathcal{D}_n$.  Strict
subordination $\mathcal{D}_n \leq \lambda\mathcal{C}_n$ with
$\lambda < 1$ therefore bounds each step's certified loss by a
proper fraction of the certified quantity: at no admissible step can
the obstruction overtake or annihilate the certified structure, and
the comparison $\mathcal{D}_n / \mathcal{C}_n \leq \lambda < 1$
is maintained along the chain.  This is persistence at the
comparison level.  (Durability of persistence on the late tail ---
uniform control of the accumulated transfer factors --- is the
separate content of SCT.3 and is not used here.) \qed
\end{proof}

\begin{theorem}[\status{C}]\label{thm:SCT3}
\textup{(SCT.3 --- Stabilized Transfer Factors Give Durable
Persistence.)}
Assume SCT.1--2.  If the effective transfer factors are
\emph{stabilized at or below unity},
$\Theta_n \leq \Theta^* \leq 1$, and strict obstruction
subordination holds with $\lambda < 1$, then
$(1+\lambda)\Theta^* \leq 1+\lambda < 2$ and persistence is
durable on the late tail.  (The branch instances satisfy the
unity bound by declaration: $\Theta_n = 1$ in the YM regime,
$\Theta_n < 1$ in the NS quasi-decay regime, and the GR import
operates at transport factor below one.)
\end{theorem}

\begin{proof}
The stepwise bound $\mathcal{C}_{n+1} \leq (1+\lambda)\Theta^*
\mathcal{C}_n$ with $(1+\lambda)\Theta^* < 2$ prevents accumulating
divergence.  Tail durability follows by induction. \qed
\end{proof}

\begin{theorem}[\status{C}]\label{thm:SCT4}
\textup{(SCT.4 --- Non-Saturation Hypothesis Implies Strict
Separation.)}
Assume SCT.1--3 and bounded obstruction $\mathcal{D}_n \leq
\mathcal{C}_n$.  Assume the \emph{non-saturation hypothesis}: there
is no admissible tail along which $\mathcal{D}_n / \mathcal{C}_n \to
1$.  Then there exist $\varepsilon > 0$ and $N$ such that
$\mathcal{D}_n \leq (1-\varepsilon)\mathcal{C}_n$ for all $n \geq N$.
\end{theorem}

\begin{proof}
If strict separation failed, then $\mathcal{D}_n/\mathcal{C}_n \to 1$
along some subsequence, contradicting the non-saturation hypothesis. \qed
\end{proof}

\begin{theorem}[\status{C}]\label{thm:SCT5}
\textup{(SCT.5 --- Irreversibility Implies Non-Saturation.)}
Assume SCT.1--4.  Assume the \emph{irreversibility condition}~(IRR):
there is no admissible tail along which
$\mathcal{D}_{n+1} - \mathcal{D}_n = \mathcal{C}_{n+1} - \mathcal{C}_n$
asymptotically (defect cannot track transport increments exactly
forever).  Then non-saturation holds.
\end{theorem}

\begin{proof}
If $\mathcal{D}_n/\mathcal{C}_n \to 1$, then asymptotically
$\Delta\mathcal{D}_n \approx \Delta\mathcal{C}_n$,
violating~(IRR).  Hence non-saturation holds. \qed
\end{proof}

\begin{theorem}[\status{C}]\label{thm:SCT6}
\textup{(SCT.6 --- Strict Dissipation Bias Implies Strict
Separation.)}
Assume SCT.1--3.  Assume the \emph{dissipation bias condition}~(DB):
there exists a branch-visible function $\Phi_n > 0$ and $\Psi_n$ with
$\mathcal{D}_{n+1} \leq \mathcal{D}_n - \Phi_n + \Psi_n$ and
$\limsup_n \Psi_n/\Phi_n < 1$, together with the lower comparison
$\Phi_n \geq c\,\mathcal{D}_n$ for some $c > 0$ on the late tail.
Then (with tail positivity of $\mathcal{C}_n$, supplied by
SCT.2--3) there exist $\varepsilon > 0$ and $N$ such that
$\mathcal{D}_n \leq (1-\varepsilon)\mathcal{C}_n$ for all $n \geq N$.
\end{theorem}

\begin{proof}
From~(DB) there is $\delta > 0$ with $\Psi_n \leq (1-\delta)\Phi_n$
for large $n$, so $\mathcal{D}_{n+1} \leq \mathcal{D}_n -
\delta\Phi_n$.  Defect experiences net loss, preventing
asymptotic saturation.  Combined with tail positivity of
$\mathcal{C}_n$ and comparison of $\Phi_n$ with $\mathcal{D}_n$,
strict separation follows. \qed
\end{proof}

\begin{theorem}[\status{C}]\label{thm:SCT6b}
\textup{(SCT.6b --- Primitive-Derived Dissipation Bias.)}
Assume SCT.1--3 and the following three primitive-derived hypotheses,
each grounded in existing spine structure:
\begin{enumerate}[label=\textup{(H\arabic*)}]
\item \textup{Quotient-loss on defect} (Books~I, III): boundary or
  interface mismatch is non-persistent; quotienting erases a fraction
  per step, contributing $\Phi_n^{\mathrm{quot}} \geq c_1
  \mathcal{D}_n$, $c_1 > 0$.
\item \textup{Finite-collar mixing loss} (Book~II, III): finite
  stabilized collar alphabet forces mixing of incompatible patterns,
  contributing $\Phi_n^{\mathrm{mix}} \geq c_2 \mathcal{D}_n$,
  $c_2 > 0$.
\item \textup{PASS-incoherent renewal} (Book~II): fresh renewal
  cannot remain phase-locked with the protected backbone, so
  $\Psi_n \leq \eta(\Phi_n^{\mathrm{quot}} + \Phi_n^{\mathrm{mix}})$
  for some $\eta < 1$.
\end{enumerate}
Setting $\Phi_n := \Phi_n^{\mathrm{quot}} + \Phi_n^{\mathrm{mix}}$,
the system satisfies strict dissipation bias:
$\mathcal{D}_{n+1} \leq \mathcal{D}_n - (1-\eta)\Phi_n$ for all
sufficiently large $n$.  In particular, no admissible tail admits
perfect long-term balance between transport and defect, and
(given tail positivity of $\mathcal{C}_n$ and lower comparison of
$\Phi_n$ against $\mathcal{D}_n$) strict separation
$\mathcal{D}_n \leq (1-\varepsilon)\mathcal{C}_n$ holds eventually.
\end{theorem}

\begin{proof}
Substituting (H3) into the loss equation gives net loss term
$(1-\eta)\Phi_n > 0$ per step.  This excludes perfect
transport-defect equilibrium.  The upgraded SCT.6 then yields strict
separation.  The canonical provenance of each hypothesis:
(H1) from Book~I quotient discipline applied to mismatch, recorded
in Book~III as transport obstruction; (H2) from Book~II finite
stabilized collar control; (H3) from Book~II PASS protection of the
backbone against halo-driven coherent coupling.  Book~VI limits the
conclusion to the declared interface regime; Book~VII prevents
promotion beyond the stated hypotheses. \qed
\end{proof}

\begin{remark}[\status{R}]
SCT.6b identifies the core asymmetry: \emph{defect is lossy,
transport is persistent}.  Defect (boundary/interface mismatch) is
non-persistent because the quotient deletes non-invariant detail
(H1) and finite collar structure prevents indefinite coherent
preservation (H2).  Renewal cannot keep pace because PASS
incoherence prevents phase-locking (H3).  The canonical summary:
\emph{the system cannot regenerate what the quotient deletes.}

The NS Branch uses SCT.6b to upgrade bounded safe-tail control to
strict positive reserve (NS.SB2).  The YM Branch satisfies a
stronger condition ($\mathcal{D}/\mathcal{C} \to 0$), which
automatically implies the SCT.6b hypotheses.
\end{remark}

%% ---------------------------------------------------------------
\section{Handoff to Book III}
%% ---------------------------------------------------------------

Book~III begins when stabilized local structure is no longer treated
as a purely static local fact, but as something that may persist,
transfer, or fail under admissible comparison across different contexts
and refinement stages.  The defect-bookkeeping objects of Book~I and
the boundary-control objects of Book~II are precisely the ingredients
that Book~III's transport and persistence machinery will act on.

The transition from Book~II to Book~III proceeds in stages: local
boundary law, then stabilized local structure, then admissible
comparison across contexts, and finally the persistence and transport
structure that is Book~III's subject.  Book~II supplies the first two
stages; Book~III supplies the remainder.

%% ---------------------------------------------------------------
\section{Dependency Ledger}
%% ---------------------------------------------------------------

\renewcommand{\arraystretch}{1.4}
\noindent
\small
\begin{tabular}{@{}p{3.6cm}p{0.8cm}p{3.8cm}p{4.8cm}@{}}
\hline
\textbf{Theorem} & \textbf{St.} & \textbf{Depends on} & \textbf{Exports to} \\
\hline
\ref{lem:toggle-orbit}\ Orbit spanning
  & [U] & Def.~\ref{def:TC}
  & \ref{thm:ls2-from-conditions} \\

\ref{lem:LAR-separation}\ LAR separation
  & [C] & Def.~\ref{def:LAR}, Bk.~I defs.
  & \ref{thm:ls2-from-conditions} \\

\ref{lem:DW-propagation}\ DW propagation
  & [C] & Def.~\ref{def:DW}, Bk.~I defs.
  & \ref{thm:ls2-from-conditions} \\

\ref{thm:ls2-from-conditions}\ Stab.\ Collar Det.\ Thm.
  & [C] & \ref{lem:toggle-orbit}--\ref{lem:DW-propagation}
  & \ref{cor:ls2-universal-K07}, \ref{thm:boundary-det} \\

\ref{thm:boundary-det}\ Boundary Determinacy
  & [C] & LS-2, Bk.~I stab.
  & \ref{cor:boundary-law-map}, \ref{prop:injection} \\

\ref{prop:injection}\ Injectivity
  & [C] & \ref{thm:boundary-det}
  & \ref{thm:outcome-count}, \ref{thm:governing} \\

\ref{thm:outcome-count}\ Outcome-count bound
  & [C] & \ref{prop:injection}
  & \ref{thm:boundary-capacity} \\

\ref{prop:collar-alphabet-bound}\ Collar bound
  & [C] & Def.~\ref{def:collar-alphabet}
  & \ref{thm:boundary-capacity} \\

\ref{thm:boundary-capacity}\ Boundary-Capacity Bound
  & [C] & \ref{thm:outcome-count}, \ref{prop:collar-alphabet-bound}
  & \ref{cor:log-interface}, \ref{thm:packing-contradiction} \\

\ref{thm:packing-contradiction}\ Local Exclusion Thm.
  & [C] & \ref{thm:boundary-capacity}, Phase-A
  & \ref{cor:log-contradiction}, \ref{cor:local-exclusion} \\

\ref{thm:governing}\ Boundary Stabilization Thm.
  & [C] & All of Bk.~II
  & Books~III--VII; all branch books \\

\ref{cor:local-exclusion}\ Cap.-Exceeding Exclusion
  & [C] & \ref{thm:packing-contradiction}
  & Branch books needing local control \\

\ref{lem:minimal-irreversibility}\ Min.\ Irreversibility
  & [U] & Def.\ of compression
  & \ref{thm:binary-rigidity} \\

\ref{thm:binary-rigidity}\ Binary Rigidity
  & [U] & Toggle group; \ref{lem:minimal-irreversibility}
  & \ref{thm:bit-budget}; Branch books \\

\ref{thm:binary-stability}\ Binary Stability
  & [U] & WL-1 regime; No-Asym.-Splitters
  & \ref{thm:bit-budget}; Branch books \\

\ref{thm:bit-budget}\ Bit-Budget Bound
  & [C] & \ref{lem:ls2-selectivity}; LS-2; E${}_8$ cap
  & YM Branch Book \\

\ref{thm:OC-gate}\ OC Gate: $\tau=0$
  & [U] & Def.~\ref{def:OC-move}
  & \ref{thm:local-nuclear} \\

\ref{thm:local-nuclear}\ Local--Nuclear Dichotomy
  & [U] & \ref{thm:OC-gate}; Weyl ineq.
  & YM + NS Branch Books \\

\ref{thm:PASS}\ PASS (derived)
  & [U] & CG-1--6; \ref{thm:outcome-count}; \ref{thm:boundary-capacity}
  & Books~III--VII; all branch books \\

\ref{cor:ceiling-finite}\ Ceiling Finiteness
  & [U] & Bk.~I Thm.~I.5.2 (Entropy Monotonicity)
  & \ref{cor:Epm-necessity}; YM Branch Book \\

\ref{cor:Epm-necessity}\ $E_\pm$ Countermodel
  & [U] & \ref{cor:ceiling-finite}; toggle calc.
  & O-KERNEL.O${}_2$; YM Branch Book \\
\hline
\end{tabular}
\normalsize
\renewcommand{\arraystretch}{1}

\medskip

\begin{scholium}[\status{R}]
Every result in Book~II whose proof invokes LS-2 is explicitly marked
conditional~[C].  The three lemmas of Section~\ref{sec:ls2-derived},
Theorem~\ref{thm:PASS} (PASS), Theorems~\ref{thm:binary-rigidity}
and~\ref{thm:binary-stability} (Binary Rigidity and Stability),
Theorem~\ref{thm:OC-gate} (Operator Collapse Gate),
Theorem~\ref{thm:local-nuclear} (Local--Nuclear Dichotomy), and
Corollary~\ref{cor:ceiling-finite} (Finiteness of Refinement Ceiling)
are the unconditional results proved here.  The status discipline of
Book~II is: combinatorial local structure and fiber characterization are
unconditional; boundary control is conditional on LS-2; exclusion is
conditional on LS-2 and Phase-A.  No result has a stronger status than
its declared dependencies allow.
\end{scholium}

\end{document}
